\chapter{The Thermal Instability of a Layer of Fluid Heated from Below}
\label{ch:3}

\section*{2. THE EFFECT OF ROTATION}

\section{Introduction}\label{sec:3-19}
In this chapter we shall investigate the effect of rotation on the simple problem of thermal instability considered in the last chapter. It will appear that rotation introduces a number of new elements into the problem; and some of its consequences are, at first sight, unexpected: the role of viscosity is, for example, inverted. The origin of this and other consequences of rotation can be traced to certain general theorems, relating to vorticity, in the dynamics of rotating fluids. For this reason, we shall preface the investigation of the stability by a discussion of these general theorems.

\section{The theorems of Helmholtz and Kelvin}\label{sec:3-20}
Consider an incompressible, inviscid, fluid. The equations governing it are:\footnote{The principal theorems to be established are not restricted to incompressible fluids; they apply to compressible fluids, also, if $p$ is a function of $\rho$ only. In the latter case $(\operatorname{grad} p)/\rho$ can be written as $\operatorname{grad} \chi$ where $\chi = \int dp/\rho$; and the quantity on the left-hand side of equation~\eqref{eq:3-1} continues to be the gradient of a scalar function. It is on this last circumstance that the principal theorems to be established depend; we shall not, in particular, have occasion to use the equation of continuity~\eqref{eq:3-2}.}
\begin{equation}\label{eq:3-1}
\frac{\partial u_i}{\partial t} + u_j \frac{\partial u_i}{\partial x_j} = -\frac{\partial}{\partial x_i}\!\left(\frac{p}{\rho} + V\right),
\end{equation}
and\dag
\begin{equation}\label{eq:3-2}
\frac{\partial u_i}{\partial x_i} = 0,
\end{equation}
where it has been assumed that the external forces are derivable from a potential $V$. The vorticity $\boldsymbol{\omega}$ is, of course, defined by
\begin{equation}\label{eq:3-3}
\boldsymbol{\omega} = \operatorname{curl} \mathbf{u} \quad \text{or} \quad \omega_i = \epsilon_{ijk}\frac{\partial u_k}{\partial x_j}.
\end{equation}

A curve drawn from point to point in the fluid, so that its direction is everywhere that of the instantaneous direction of $\boldsymbol{\omega}$, is called a \emph{vortex line}. The differential equations determining it are
\begin{equation}\label{eq:3-4}
\frac{dx_1}{\omega_1} = \frac{dx_2}{\omega_2} = \frac{dx_3}{\omega_3}.
\end{equation}

If through every point of a small closed curve we draw the corresponding vortex line, we obtain a tube which is called a \emph{vortex tube}.

Since
\begin{equation}\label{eq:3-5}
\operatorname{div} \boldsymbol{\omega} = 0,
\end{equation}
it follows from Gauss's theorem that the integral of the normal component of $\boldsymbol{\omega}$ over any closed surface is zero:
\begin{equation}\label{eq:3-6}
\oint_S \boldsymbol{\omega} \cdot d\mathbf{S} = 0 \quad (S\ \text{any closed surface}).
\end{equation}

Apply this result to the closed surface bounding an element of fluid contained in a section of the vortex tube (see Fig.~18). The sides of the vortex tube clearly do not contribute to the surface integral and equation~\eqref{eq:3-6}, in this case, gives
\begin{equation}\label{eq:3-7}
\int_{S_1} \boldsymbol{\omega} \cdot d\mathbf{S}_1 = \int_{S_2} \boldsymbol{\omega} \cdot d\mathbf{S}_2.
\end{equation}

\figurePlaceholder{18}{A vortex tube.}

\emph{The flux of vorticity across any section of a vortex tube is the same}; it, therefore, represents a characteristic of the tube. If we consider a vortex tube of infinitesimal cross-section, the flux across a section $d\mathbf{S}$ normal to itself is $|\boldsymbol{\omega}|dS$. By the theorem we have just proved, $|\boldsymbol{\omega}|dS$ is constant along the tube; it cannot therefore terminate at any point in the fluid. \emph{Vortex lines must either be closed, or terminate on the boundaries.}

We shall now obtain an equation of motion for the vorticity. First, we observe that
\begin{equation}\label{eq:3-8}
\epsilon_{ijk}\, u_j\, \omega_k = \epsilon_{ijk}\,\epsilon_{klm}\, u_j \frac{\partial u_m}{\partial x_l}
= (\delta_{il}\,\delta_{jm} - \delta_{im}\,\delta_{jl})\, u_j \frac{\partial u_m}{\partial x_l}
= u_j\frac{\partial u_j}{\partial x_i} - u_j\frac{\partial u_i}{\partial x_j},
\end{equation}
or
\begin{equation}\label{eq:3-9}
u_j\frac{\partial u_i}{\partial x_j} = \frac{1}{2}\frac{\partial}{\partial x_i}|\mathbf{u}|^2 - \epsilon_{ijk}\, u_j\, \omega_k.
\end{equation}

Using this relation in~\eqref{eq:3-1}, we have
\begin{equation}\label{eq:3-10}
\frac{\partial u_i}{\partial t} - \epsilon_{ijk}\, u_j\, \omega_k = -\frac{\partial}{\partial x_i}\!\left(\frac{p}{\rho} + \tfrac{1}{2}|\mathbf{u}|^2 + V\right).
\end{equation}

Letting
\begin{equation}\label{eq:3-11}
\varpi = p/\rho + \tfrac{1}{2}|\mathbf{u}|^2 + V,
\end{equation}
we can write
\begin{equation}\label{eq:3-12}
\frac{\partial \mathbf{u}}{\partial t} - \mathbf{u}\times\boldsymbol{\omega} = \operatorname{grad}\varpi.
\end{equation}

Taking the curl of this equation, we obtain
\begin{equation}\label{eq:3-13}
\frac{\partial\boldsymbol{\omega}}{\partial t} - \operatorname{curl}(\mathbf{u}\times\boldsymbol{\omega}) = 0;
\end{equation}
this is the required equation of motion for $\boldsymbol{\omega}$. The principal theorem on vorticity follows from this equation.

Consider a surface $S$ enclosed by a simple closed contour $C$. Let $d\mathbf{S}$ be an element of this surface. Multiplying equation~\eqref{eq:3-13} scalarly by $d\mathbf{S}$, and integrating over $S$, we obtain
\begin{equation}\label{eq:3-14}
\int_S \frac{\partial\boldsymbol{\omega}}{\partial t}\cdot d\mathbf{S} - \int_S \operatorname{curl}(\mathbf{u}\times\boldsymbol{\omega})\cdot d\mathbf{S} = 0.
\end{equation}

\figurePlaceholder{19}{Illustrating the Helmholtz--Kelvin theorem.}

Transforming the second of these integrals by Stokes's theorem, we have
\begin{equation}\label{eq:3-15}
\int_S \frac{\partial\boldsymbol{\omega}}{\partial t}\cdot d\mathbf{S} + \oint_C \boldsymbol{\omega}\cdot(\mathbf{u}\times d\mathbf{s}) = 0,
\end{equation}
where $d\mathbf{s}$ is an element of arc of the contour $C$. The physical meaning of this last equation becomes clearer if we multiply the equation by an infinitesimal interval of time $\Delta t$, and then interpret the different terms. We have
\begin{equation}\label{eq:3-16}
\Delta t \int_S \frac{\partial\boldsymbol{\omega}_t}{\partial t}\cdot d\mathbf{S} + \oint_{C_t} \boldsymbol{\omega}_t\cdot(\mathbf{u}\,\Delta t \times d\mathbf{s}) = 0,
\end{equation}
where we have inserted a subscript $t$ to the various quantities to emphasize that they are the instantaneous values at the time indicated (see Fig.~19). Now $\mathbf{u}\,\Delta t \times d\mathbf{s}$ is the area swept out by the fluid elements along $d\mathbf{s}$ in a time $\Delta t$. Accordingly, the integral along $C_t$ in~\eqref{eq:3-16} is, in reality, a \emph{surface integral} over the elementary strip connecting the contour $C_t$ with the contour $C_{t+\Delta t}$ to which $C_t$ will be carried in the time $\Delta t$. Applying the result~\eqref{eq:3-6} to the closed surface consisting of $S_t$ (bounded by $C_t$), $S_{t+\Delta t}$ (bounded by $C_{t+\Delta t}$), and the strip connecting $C_t$ and $C_{t+\Delta t}$, we have\footnote{Note that the validity of this equation depends on $\operatorname{div}\boldsymbol{\omega} = 0$.}
\begin{equation}\label{eq:3-17}
\int_{S_{t+\Delta t}} \boldsymbol{\omega}_t\cdot(\mathbf{u}\,\Delta t \times d\mathbf{s}) = \int_{S_{t+\Delta t}} \boldsymbol{\omega}_t\cdot d\mathbf{S} - \int_{S_t} \boldsymbol{\omega}_t\cdot d\mathbf{S}.
\end{equation}

Using this in equation~\eqref{eq:3-16}, we obtain
\begin{equation}\label{eq:3-18}
\Delta t \int_{S_t} \frac{\partial\boldsymbol{\omega}_t}{\partial t}\cdot d\mathbf{S} + \int_{S_{t+\Delta t}} \boldsymbol{\omega}_t\cdot d\mathbf{S} = \int_{S_t} \boldsymbol{\omega}_t\cdot d\mathbf{S}.
\end{equation}

With an error of order $(\Delta t)^2$, we can replace the integral over $S_t$ on the left-hand side by an integral over $S_{t+\Delta t}$; we thus find,
\begin{equation}\label{eq:3-19}
\int_{S_{t+\Delta t}} \!\left(\boldsymbol{\omega}_t + \Delta t\frac{\partial\boldsymbol{\omega}_t}{\partial t}\right)\!\cdot d\mathbf{S} = \int_{S_t} \boldsymbol{\omega}_t\cdot d\mathbf{S} + O(\Delta t^2),
\end{equation}
or
\begin{equation}\label{eq:3-20}
\int_{S_{t+\Delta t}} \boldsymbol{\omega}_{t+\Delta t}\cdot d\mathbf{S} = \int_{S_t} \boldsymbol{\omega}_t\cdot d\mathbf{S} + O(\Delta t^2).
\end{equation}

Passing to the limit $\Delta t = 0$, we obtain
\begin{equation}\label{eq:3-21}
\frac{d}{dt}\int_{S_t}\boldsymbol{\omega}\cdot d\mathbf{S} = 0.
\end{equation}

We have thus proved: \emph{The integral of the normal component of\/ $\boldsymbol{\omega}$ over any surface $S$ bounded by a closed curve remains constant as we follow the surface $S$ with the motion of the fluid elements constituting it}:
\begin{equation}\label{eq:3-22}
\int_{S_t}\boldsymbol{\omega}_t\cdot d\mathbf{S} = \text{constant}.
\end{equation}

This is the principal theorem on vorticity due to Helmholtz and Kelvin. We may state it also in the form: \emph{the strength of a vortex tube is an integral of the equations of motion}.

Transforming~\eqref{eq:3-22} by Stokes's theorem, we have
\begin{equation}\label{eq:3-23}
\int_S \boldsymbol{\omega}\cdot d\mathbf{S} = \int_S \operatorname{curl}\mathbf{u}\cdot d\mathbf{S} \stackrel{*}{=} \oint_C \mathbf{u}\cdot d\mathbf{s} = \text{constant}.
\end{equation}

The integral of $\mathbf{u}$ along any closed curve is the \emph{circulation} along the curve. By~\eqref{eq:3-23}, the circulation along any closed curve $C$ remains constant as we follow the curve $C$ with the motion of the fluid elements constituting it.

It is customary to derive the constancy of the circulation along a closed curve first, and then deduce from it the constancy of the flux of the vorticity across a surface. We have followed a reverse procedure to make comparisons with the motions of the magnetic lines of force in hydromagnetics (Chapter~IV, \S\,38\,(b)).

A further important fact is that vortex lines retain their character as vortex lines as we follow the fluid motion. This follows from equation~\eqref{eq:3-13} written somewhat differently. Since
\begin{equation}\label{eq:3-24}
\epsilon_{ijk}\,\epsilon_{klm}\, u_j\, \omega_m = (\delta_{il}\,\delta_{jm} - \delta_{im}\,\delta_{jl})\frac{\partial}{\partial x_j}(u_j\,\omega_i - u_i\,\omega_j) = \frac{\partial}{\partial x_j}(u_j\,\omega_i - u_i\,\omega_j),
\end{equation}
we can write
\begin{equation}\label{eq:3-25}
\frac{\partial\omega_i}{\partial t} + \frac{\partial}{\partial x_j}(u_j\,\omega_i - u_i\,\omega_j) = 0.
\end{equation}

Using the solenoidal character\dag\ of $\mathbf{u}$ and $\boldsymbol{\omega}$, we can rewrite this last equation in the form
\begin{equation}\label{eq:3-26}
\frac{\partial\omega_i}{\partial t} + u_j\frac{\partial\omega_i}{\partial x_j} = \omega_j\frac{\partial u_i}{\partial x_j},
\end{equation}
or
\begin{equation}\label{eq:3-27}
\frac{\partial\boldsymbol{\omega}}{\partial t} + (\mathbf{u}\cdot\operatorname{grad})\boldsymbol{\omega} = \frac{d\boldsymbol{\omega}}{dt} = (\boldsymbol{\omega}\cdot\operatorname{grad})\mathbf{u}.
\end{equation}

Consider now a vortex line. Let $A$ and $B$ be two adjacent points on it. Since the element of arc joining $A$ and $B$ is in the direction of $\boldsymbol{\omega}$,
\begin{equation}\label{eq:3-28}
\overrightarrow{AB} = \delta q\,\boldsymbol{\omega},
\end{equation}
where $\delta q$ is an infinitesimal constant. The velocities at $A$ and $B$ differ by the amount
\begin{equation}\label{eq:3-29}
\mathbf{u}_B - \mathbf{u}_A = (\overrightarrow{AB}\cdot\operatorname{grad})\mathbf{u}_A = \delta q\,(\boldsymbol{\omega}\cdot\operatorname{grad})\mathbf{u}_A;
\end{equation}
or, making use of equation~\eqref{eq:3-27}, we have
\begin{equation}\label{eq:3-30}
\mathbf{u}_B - \mathbf{u}_A = \delta q\,\frac{d\boldsymbol{\omega}}{dt}.
\end{equation}

During a time $\Delta t$, the points $A$ and $B$ will suffer displacements of amounts $\mathbf{u}_A\,\Delta t$ and $\mathbf{u}_B\,\Delta t$, respectively. If the displaced positions of $A$ and $B$ are $A'$ and $B'$, then
\begin{equation}\label{eq:3-31}
\overrightarrow{A'B'} = \overrightarrow{AB} + (\mathbf{u}_B - \mathbf{u}_A)\Delta t
= \delta q\!\left(\boldsymbol{\omega} + \frac{d\boldsymbol{\omega}}{dt}\,\Delta t\right) = \delta q\,\boldsymbol{\omega}_{t+\Delta t}.
\end{equation}

In other words, $\overrightarrow{A'B'}$ is tangential to the vortex line through $A'$ at time $t+\Delta t$. \emph{A vortex line, therefore, consists of the same fluid elements: it moves with the fluid like a material substance.} We may if we like say: \emph{the vortex lines are permanently attached to the fluid.} If a vortex line should disappear in a real fluid, it can only be because of dissipation by viscosity.

\section{The equations of hydrodynamics in a rotating frame of reference}\label{sec:3-21}
Consider a fluid in rotation about some fixed axis with a constant angular velocity $\boldsymbol{\Omega}$. It will be convenient to describe the motions which occur in it as they will appear to an observer at rest in a frame rotating about the same axis and with the same angular velocity. In such a rotating frame of reference, quantities which will be recognized as velocities and accelerations by an observer at rest in it are not the same as the velocities and accelerations as recognized by an observer at rest with respect to the fixed inertial frame $(\xi, \eta, \zeta)$ with respect to which the chosen coordinate system $(x,y,z)$ is rotating with the angular velocity $\Omega$ about the $\zeta$-axis. Also, the $z$- and the $\zeta$-axes will be assumed to coincide. The transformation relating the two coordinate systems is
\begin{align}
x &= +\xi\cos\Omega t + \eta\sin\Omega t, \label{eq:3-32}\\
y &= -\xi\sin\Omega t + \eta\cos\Omega t, \notag\\
z &= +\zeta. \notag
\end{align}

A vector $\mathbf{q}$, with components $q_\xi$, $q_\eta$, and $q_\zeta$ in the inertial frame, will have the following components along the instantaneous directions of the axes of the rotating frame:
\begin{align}
q_x^{(0)} &= +q_\xi\cos\Omega t + q_\eta\sin\Omega t, \label{eq:3-33}\\
q_y^{(0)} &= -q_\xi\sin\Omega t + q_\eta\cos\Omega t, \notag\\
q_z^{(0)} &= +q_\zeta. \notag
\end{align}

We have distinguished the components of the vector obtained by this transformation by a superscript (0) because they are not necessarily the same quantities which will be recognized by the observer, at rest in the rotating frame, as having the meaning of $\mathbf{q}$. The superscript (0) will distinguish the quantities which will have the same \emph{absolute}, as distinct from \emph{relative}, meanings. The reason for this distinction will become clear when we consider the quantities which have the meanings of velocity and acceleration in the rotating frame and their relations to the `absolute' velocity and acceleration.

By differentiating equation~\eqref{eq:3-32} with respect to $t$, we obtain
\begin{align}
\frac{dx}{dt} &= \left(+\frac{d\xi}{dt}\cos\Omega t + \frac{d\eta}{dt}\sin\Omega t\right) - \Omega(\xi\sin\Omega t - \eta\cos\Omega t), \label{eq:3-34}\\
\frac{dy}{dt} &= \left(-\frac{d\xi}{dt}\sin\Omega t + \frac{d\eta}{dt}\cos\Omega t\right) - \Omega(\xi\cos\Omega t + \eta\sin\Omega t). \label{eq:3-35}
\end{align}

Clearly, $dx/dt$ and $dy/dt$ are the quantities which will be recognized by the observer at rest in the rotating frame as the components $u_x$ and $u_y$ of the velocities of the fluid element along the $x$- and the $y$-directions; whereas the quantities grouped together in the first brackets on the right-hand sides of equations~\eqref{eq:3-34} and~\eqref{eq:3-35} are the components along the instantaneous directions of $x$ and $y$ of what is the velocity in the inertial frame; they are the quantities which must be distinguished by the superscript (0). Equations~\eqref{eq:3-34} and~\eqref{eq:3-35} relate $u_x$ and $u_y$ with $u_x^{(0)}$ and $u_y^{(0)}$; we have
\begin{equation}\label{eq:3-36}
u_x = u_x^{(0)} + \Omega y; \quad u_y = u_y^{(0)} - \Omega x; \quad u_z = u_z^{(0)}.
\end{equation}

These equations can be combined into the single vector equation
\begin{equation}\label{eq:3-37}
\mathbf{u} = \mathbf{u}^{(0)} - \boldsymbol{\Omega}\times\mathbf{r}.
\end{equation}

Now differentiating equations~\eqref{eq:3-34} and~\eqref{eq:3-35} once again with respect to $t$, we obtain
\begin{align}
\frac{d^2 x}{dt^2} &= \left(+\frac{d^2\xi}{dt^2}\cos\Omega t + \frac{d^2\eta}{dt^2}\sin\Omega t\right) + 2\Omega\!\left(-\frac{d\xi}{dt}\sin\Omega t + \frac{d\eta}{dt}\cos\Omega t\right) - \Omega^2 x, \label{eq:3-38}\\
\frac{d^2 y}{dt^2} &= \left(-\frac{d^2\xi}{dt^2}\sin\Omega t + \frac{d^2\eta}{dt^2}\cos\Omega t\right) + 2\Omega\!\left(-\frac{d\xi}{dt}\cos\Omega t - \frac{d\eta}{dt}\sin\Omega t\right) - \Omega^2 y. \notag
\end{align}

The equations are equivalent to
\begin{align}
\frac{d^2 x}{dt^2} &= \left(\frac{du_x^{(0)}}{dt}\right)^{(0)} + 2\Omega u_y^{(0)} - \Omega^2 x, \label{eq:3-39}\\
\frac{d^2 y}{dt^2} &= \left(\frac{du_y^{(0)}}{dt}\right)^{(0)} - 2\Omega u_x^{(0)} - \Omega^2 y. \notag
\end{align}

Substituting for $u_x^{(0)}$ and $u_y^{(0)}$ from~\eqref{eq:3-36} in the foregoing equations, we obtain
\begin{align}
\left(\frac{du_x^{(0)}}{dt}\right)^{(0)} &= \frac{du_x}{dt} - 2\Omega u_y - \Omega^2 x, \label{eq:3-40}\\
\left(\frac{du_y^{(0)}}{dt}\right)^{(0)} &= \frac{du_y}{dt} + 2\Omega u_x - \Omega^2 y, \notag\\
\left(\frac{du_z^{(0)}}{dt}\right)^{(0)} &= \frac{du_z}{dt}. \notag
\end{align}

These equations can be combined into the single vector equation
\begin{equation}\label{eq:3-41}
\left(\frac{d\mathbf{u}^{(0)}}{dt}\right)^{(0)} = \frac{d\mathbf{u}}{dt} + 2\boldsymbol{\Omega}\times\mathbf{u} - \tfrac{1}{2}\operatorname{grad}(|\boldsymbol{\Omega}\times\mathbf{r}|^2).
\end{equation}

The term $2\boldsymbol{\Omega}\times\mathbf{u}$ in this equation represents the Coriolis acceleration and the term $-\frac{1}{2}\operatorname{grad}(|\boldsymbol{\Omega}\times\mathbf{r}|^2)$ represents the centrifugal force.

For differentiations with respect to the space coordinates, there are no complicating considerations such as we have just explained for differentiations with respect to time leading to velocities and accelerations. Accordingly, the standard hydrodynamical equation (equation~II~(17)),
\begin{equation}\label{eq:3-42}
\frac{du_i}{dt} = X_i + \frac{1}{\rho}\frac{\partial P_{ij}}{\partial x_j}
\end{equation}
(where $X_i$ is the $i$th component of whatever external forces may be acting on the fluid and $P_{ij}$ is the stress tensor), in the rotating frame of reference, becomes
\begin{equation}\label{eq:3-43}
\frac{du_i}{dt} + u_j\frac{\partial u_i}{\partial x_j} = X_i + \frac{1}{\rho}\frac{\partial P_{ij}}{\partial x_j} + \frac{\partial}{\partial x_i}\!\left(\tfrac{1}{2}|\boldsymbol{\Omega}\times\mathbf{r}|^2\right) + 2\epsilon_{ijk}\, u_j\,\Omega_k.
\end{equation}

The equation of continuity and the equation of heat conduction remain unaffected.

For an incompressible fluid, the equations of motion take the explicit form
\begin{equation}\label{eq:3-44}
\frac{\partial u_i}{\partial t} + u_j\frac{\partial u_i}{\partial x_j} = X_i - \frac{\partial}{\partial x_i}\!\left(\frac{p}{\rho} - \tfrac{1}{2}|\boldsymbol{\Omega}\times\mathbf{r}|^2\right) + \nu\nabla^2 u_i + 2\epsilon_{ijk}\, u_j\,\Omega_k.
\end{equation}

The equation of continuity in the form~\eqref{eq:3-2} continues to hold.

\section{The Taylor--Proudman theorem}\label{sec:3-22}
Consider the equation of motion~\eqref{eq:3-44} for an inviscid fluid in case the external forces are derivable from a potential $V$. We then have
\begin{equation}\label{eq:3-45}
\frac{\partial u_i}{\partial t} + u_j\frac{\partial u_i}{\partial x_j} = -\frac{\partial}{\partial x_i}\!\left(\frac{p}{\rho} - \tfrac{1}{2}|\boldsymbol{\Omega}\times\mathbf{r}|^2 + V\right) + 2\epsilon_{ijk}\, u_j\,\Omega_k.
\end{equation}

Making use of the identity~\eqref{eq:3-9}, we can write
\begin{equation}\label{eq:3-46}
\frac{\partial u_i}{\partial t} - \epsilon_{ijk}\, u_j\, \omega_k = -\frac{\partial\varpi}{\partial x_i} + 2\epsilon_{ijk}\, u_j\,\Omega_k,
\end{equation}
where now
\begin{equation}\label{eq:3-47}
\varpi = p/\rho + \tfrac{1}{2}|\mathbf{u}|^2 + V - \tfrac{1}{2}|\boldsymbol{\Omega}\times\mathbf{r}|^2.
\end{equation}

Alternatively, we can also write
\begin{equation}\label{eq:3-48}
\frac{\partial\mathbf{u}}{\partial t} - \mathbf{u}\times(\boldsymbol{\omega}+2\boldsymbol{\Omega}) = -\operatorname{grad}\varpi,
\end{equation}
and taking the curl of this equation, we have
\begin{equation}\label{eq:3-49}
\frac{\partial\boldsymbol{\omega}}{\partial t} - \operatorname{curl}[\mathbf{u}\times(\boldsymbol{\omega}+2\boldsymbol{\Omega})] = 0.
\end{equation}

Since $\boldsymbol{\Omega}$ is a constant vector, it is clear from a comparison of equations~\eqref{eq:3-13} and~\eqref{eq:3-49} that \emph{in a rotating frame of reference $\boldsymbol{\omega}+2\boldsymbol{\Omega}$ plays the role of\/ $\boldsymbol{\omega}$ in an inertial frame}. We can, therefore, conclude that now
\begin{equation}\label{eq:3-50}
\int_S (\boldsymbol{\omega}+2\boldsymbol{\Omega})\cdot d\mathbf{S} = \text{constant}
\end{equation}
as we follow a surface $S$ (enclosed by a simple closed curve $C$) with the motion in the rotating frame. An equivalent statement is:

\emph{The circulation round a closed contour $C$ bounding a surface $S$ $+$ $2\boldsymbol{\Omega}\times$(the projection of the area of $S$ on a plane perpendicular to $\boldsymbol{\Omega}$)}
\begin{equation}\label{eq:3-51}
= \text{constant}.
\end{equation}

In this form, the theorem appears to have been first stated by Taylor. When the motions are slow and steady, the theorem asserts:

\begin{equation}\label{eq:3-52}
\textit{The projection of the area of\/ $S$ on a plane perpendicular to\/ $\boldsymbol{\Omega}$ is a constant.}
\end{equation}

What this implies for the motions themselves can be inferred more directly from equation~\eqref{eq:3-49}. In a state of steady, slow motions when the squares of the velocity components can be neglected, equation~\eqref{eq:3-49} requires that
\begin{equation}\label{eq:3-53}
\operatorname{curl}(\mathbf{u}\times\boldsymbol{\Omega}) = 0,
\end{equation}
or (cf.\ equation~\eqref{eq:3-24})
\begin{equation}\label{eq:3-54}
\frac{\partial}{\partial x_j}(\omega_i\,\Omega_j - u_j\,\Omega_i) = 0.
\end{equation}

Since $\boldsymbol{\Omega}$ is a constant vector and $\mathbf{u}$ is solenoidal, equation~\eqref{eq:3-54} reduces to
\begin{equation}\label{eq:3-55}
\Omega_j\frac{\partial u_i}{\partial x_j} = 0.
\end{equation}

\emph{The motions cannot, therefore, vary in the direction of\/ $\boldsymbol{\Omega}$. In other words, all steady slow motions in a rotating inviscid fluid are necessarily two-dimensional.} This is the form in which the Taylor--Proudman theorem is generally stated; equation~\eqref{eq:3-50} is, however, a more complete statement.

What the condition~\eqref{eq:3-55} implies is this: the fact that the motions transverse to $\boldsymbol{\Omega}$ cannot vary in the direction of $\boldsymbol{\Omega}$ means that any two fluid elements which are initially on a line parallel to $\boldsymbol{\Omega}$ will always remain on that line; and the fact that motions in the direction of $\boldsymbol{\Omega}$ cannot also vary along this direction means that two fluid elements which are initially a certain distance apart (in the direction of $\boldsymbol{\Omega}$) will always remain at this distance apart. This predicted behaviour of slow, steady motions in a rotating fluid was strikingly demonstrated by Taylor as follows.

Water in a tank was first made to rotate steadily as a solid body. A small motion was then communicated to it and a few drops of coloured liquid (or ink) were inserted. The slow movement of the fluid then drew this coloured portion of the fluid into sheets; and these sheets

\figurePlaceholder{20}{Illustrating the Taylor--Proudman theorem in a rotating cylinder of water. In the experiment, ink was poured into the rotating cylinder and stirred $1\frac{1}{2}$~min later by dipping a stationary rod for 15~sec. Note the thinness of many of the filaments into which the ink is drawn in accordance with the Taylor--Proudman theorem.}

\noindent always remained parallel to the axis of rotation. Taylor noted: `The accuracy with which they remained parallel to the axis of rotation is quite extraordinary.' These experiments of Taylor have been repeated by Fultz; one of his photographs illustrating this phenomenon is reproduced in Fig.~20.

\section{The propagation of waves in a rotating fluid}\label{sec:3-23}
We shall conclude this excursion into general hydrodynamical theorems by showing that the equations of motion in a rotating fluid allow periodic solutions representing the propagation of waves.

In the absence of external forces, or in the presence of forces which are derivable from a potential, the equations of motion governing a viscous incompressible fluid take the form (cf.\ equation~\eqref{eq:3-44})
\begin{equation}\label{eq:3-56}
\frac{\partial u_i}{\partial t} + u_j\frac{\partial u_i}{\partial x_j} = -\frac{\partial\varpi}{\partial x_i} + \nu\nabla^2 u_i + 2\epsilon_{ijk}\, u_j\,\Omega_k.
\end{equation}

We seek solutions of this equation with a space-time dependence given by
\begin{equation}\label{eq:3-57}
e^{i(\mathbf{k}\cdot\mathbf{r}+pt)}.
\end{equation}

For solutions having this space-time dependence,
\begin{equation}\label{eq:3-58}
\frac{\partial}{\partial t} = ip, \quad \frac{\partial}{\partial x_j} = ik_j, \quad \text{and} \quad \nabla^2 = -k^2;
\end{equation}
and equation~\eqref{eq:3-56} gives
\begin{equation}\label{eq:3-59}
(ip + \nu k^2)u_i + ik_j\, u_j\, u_i = -ik_i\,\varpi + 2\epsilon_{ijk}\, u_j\,\Omega_k;
\end{equation}
while the equation of continuity gives
\begin{equation}\label{eq:3-60}
k_j\, u_j = 0.
\end{equation}

This last equation implies that \emph{the waves are transverse}.

In view of~\eqref{eq:3-60}, equation~\eqref{eq:3-59} becomes
\begin{equation}\label{eq:3-61}
n u_i + k_j\,\varpi + 2i\epsilon_{ijk}\, u_j\,\Omega_k = 0,
\end{equation}
where we have written
\begin{equation}\label{eq:3-62}
n = p - i\nu k^2.
\end{equation}

Multiplying equation~\eqref{eq:3-61} scalarly by $u_i$, $k_i$, and $\Omega_i$, we obtain, respectively,
\begin{equation}\label{eq:3-63}
\begin{aligned}
n u_i^2 &= 0, \\
k^2\,\varpi + 2i\epsilon_{ijk}\, k_i\, u_j\,\Omega_k &= 0, \\
n u_i\,\Omega_i + \varpi k_i\,\Omega_i &= 0.
\end{aligned}
\end{equation}

It is now convenient to choose the orientation of the coordinate system such that
\begin{equation}\label{eq:3-64}
\boldsymbol{\Omega} = (\Omega_x, 0, \Omega_z) \quad \text{and} \quad \mathbf{k} = (0,0,k);
\end{equation}
this clearly implies no loss of generality. With this choice of the coordinate system, equation~\eqref{eq:3-60} gives
\begin{equation}\label{eq:3-65}
k u_z = 0 \quad \text{or} \quad u_z = 0.
\end{equation}

Equations~\eqref{eq:3-63} now give
\begin{align}
n(u_x^2 + u_y^2) &= 0, \label{eq:3-66}\\
2\Omega_z\, u_y &= +k\varpi, \label{eq:3-67}\\
n\Omega_x\, u_x &= -k\varpi\Omega_z. \label{eq:3-68}
\end{align}

These equations enable us to express the amplitudes $u_x$ and $\varpi$ in terms of $u_z$; we find
\begin{equation}\label{eq:3-69}
u_y = +i\frac{n}{2\Omega_z}u_x; \quad \varpi = -\frac{n\Omega_x}{k\Omega_z}u_x; \quad u_z = 0.
\end{equation}

Substituting the first of these relations in equation~\eqref{eq:3-66}, we obtain
\begin{equation}\label{eq:3-70}
n\!\left(1 - \frac{n^2}{4\Omega_z^2}\right)\!u_x^2 = 0.
\end{equation}

From this equation it follows that $n$ is real; and ignoring the root $n = 0$, we have
\begin{equation}\label{eq:3-71}
n = \pm 2\Omega_z = \pm 2\Omega\cos\vartheta,
\end{equation}
where $\vartheta$ is the inclination of the direction of wave-propagation to the direction of $\boldsymbol{\Omega}$. The gyration frequency $p$ is given by (cf.\ equation~\eqref{eq:3-62})
\begin{equation}\label{eq:3-72}
p = n + i\nu k^2 = \pm 2\Omega\cos\vartheta + i\nu k^2.
\end{equation}

The waves are, therefore, damped; and the damping constant is $\nu k^2$.

With $n$ given by equation~\eqref{eq:3-72}, the expressions for the amplitudes become
\begin{equation}\label{eq:3-73}
u_x = \pm i u_y, \quad \varpi = \mp\frac{2\Omega}{k}u_x\sin\vartheta, \quad u_z = 0.
\end{equation}

\emph{The waves are, therefore, transverse, circularly polarized, and damped.}

In the limit of zero viscosity, the phase velocity is given by
\begin{equation}\label{eq:3-74}
V = \frac{p}{k} = \pm\frac{2\Omega}{k}\cos\vartheta.
\end{equation}

This is the \emph{phase velocity}. The \emph{group velocity} is given by
\begin{equation}\label{eq:3-75}
\frac{\partial p}{\partial k_i} = \pm 2\frac{\partial}{\partial k_i}\!\left(\frac{\Omega_j k_j}{k}\right) = \pm 2\!\left(\frac{\Omega_i}{k} - \frac{\Omega_j k_j}{k^3}k_i\right),
\end{equation}
or
\begin{equation}\label{eq:3-76}
\frac{\partial p}{\partial \mathbf{k}} = \pm\frac{2}{k^3}\mathbf{k}\times(\boldsymbol{\Omega}\times\mathbf{k}).
\end{equation}

Finally, it should be noted that the foregoing solutions representing the propagation of waves are not limited to infinitesimal amplitudes.

\section{The problem of thermal instability in a rotating fluid: general considerations}\label{sec:3-24}
From the theorems proved in \S\S\,22 and~23 it is apparent that rotation should have a profound effect on the onset of thermal instability. For convection implies that motions which occur have, necessarily, a three-dimensional character; this the Taylor--Proudman theorem expressly forbids for an inviscid fluid so long as the non-linear terms in the equations of motion are neglected and the motions are steady. The first of these two conditions is certainly met in a linear stability theory; and the second obtains for steady convection. Therefore, in contrast to non-rotating fluids, an inviscid fluid in rotation should be expected to be thermally stable for \emph{all} adverse temperature gradients. Indeed, only in the presence of viscosity can thermal instability arise; for only then can the Taylor--Proudman theorem be violated.

There is another factor to remember. While the Taylor--Proudman theorem forbids any variation of the velocity in the direction of $\boldsymbol{\Omega}$, it does not forbid oscillatory motions. As we have seen in \S\,23, the propagation of transverse waves is possible under the same circumstances. For this reason, we may anticipate that overstability as a means of developing instability may, under certain conditions, play a decisive role. We shall see that this is indeed the case.

\section{The perturbation equations}\label{sec:3-25}
Consider, then, an infinite horizontal layer of fluid which is kept rotating at a constant rate. Let $\boldsymbol{\Omega}$ denote the angular velocity of rotation. As we have shown explicitly in Chapter~II, \S\,9 for problems of the type we are considering, it will suffice to work with the relevant equations of motion and heat conduction in the Boussinesq approximation. The only additional factors which we have to allow for now are the effects of Coriolis acceleration and centrifugal force in the equations of motion. Thus, in view of equation~\eqref{eq:3-41}, we replace equation~II~(43) by
\begin{equation}\label{eq:3-77}
\frac{\partial u_i}{\partial t} + u_j\frac{\partial u_i}{\partial x_j} = -\frac{\partial}{\partial x_i}\!\left(\frac{\delta p}{\rho_0} - \tfrac{1}{2}|\boldsymbol{\Omega}\times\mathbf{r}|^2\right) + \left(1+\frac{\delta\rho}{\rho_0}\right)X_i + \nu\nabla^2 u_i + 2\epsilon_{ijk}\, u_j\,\Omega_k.
\end{equation}

The remaining equations (equations~II~(41), (44), and (46)) are unaffected.

We envisage now an initial state in which a steady adverse temperature gradient $\beta$ is maintained and there are no motions. The characterization of this initial state differs in no respect from that described in Chapter~II, \S\,9; and the equations governing small perturbations can be obtained in exactly the same way. Thus, equation~II~(55) is replaced by
\begin{equation}\label{eq:3-78}
\frac{\partial u_i}{\partial t} = -\frac{\partial}{\partial x_i}\!\left(\frac{\delta p}{\rho_0}\right) + g\alpha\Theta\,\lambda_i + \nu\nabla^2 u_i + 2\epsilon_{ijk}\, u_j\,\Omega_k,
\end{equation}
where $\boldsymbol{\lambda} = (0,0,1)$ is a unit vector in the direction of the vertical. We also have (equations~II~(56) and~(57))
\begin{equation}\label{eq:3-79}
\frac{\partial\Theta}{\partial t} = \beta\lambda_j u_j + \kappa\nabla^2\Theta
\end{equation}
and
\begin{equation}\label{eq:3-80}
\frac{\partial u_j}{\partial x_j} = 0.
\end{equation}

As in Chapter~II, \S\,9, we can eliminate the term in $\delta p/\rho_0$ in equation~\eqref{eq:3-78} by applying the operator $\epsilon_{ijk}\,\partial/\partial x_j$ to the $k$-component of the equation; since (cf.\ equation~\eqref{eq:3-24})
\begin{equation}\label{eq:3-81}
\epsilon_{ijk}\frac{\partial}{\partial x_j}\epsilon_{klm}\, u_l\,\Omega_m = \frac{\partial}{\partial x_j}(u_i\,\Omega_j - u_j\,\Omega_i) = \Omega_j\frac{\partial u_i}{\partial x_j},
\end{equation}
the result is (cf.\ equation~II~(67))
\begin{equation}\label{eq:3-82}
\frac{\partial\omega_i}{\partial t} = g\alpha\epsilon_{ijk}\frac{\partial\Theta}{\partial x_j}\lambda_k + \nu\nabla^2\omega_i + 2\Omega_j\frac{\partial u_i}{\partial x_j}.
\end{equation}

Taking the curl of this equation once again, we obtain (cf.\ equation~II~(72))
\begin{equation}\label{eq:3-83}
\frac{\partial}{\partial t}(\nabla^2 u_i) = g\alpha\!\left(\lambda_i\nabla^2\Theta - \lambda_j\frac{\partial^2\Theta}{\partial x_i\,\partial x_j}\right) + \nu\nabla^4 u_i - 2\Omega_j\frac{\partial\omega_i}{\partial x_j}.
\end{equation}

Now multiplying equations~\eqref{eq:3-82} and~\eqref{eq:3-83} by $\lambda_i$, we get
\begin{equation}\label{eq:3-84}
\frac{\partial\zeta}{\partial t} = \nu\nabla^2\zeta + 2\Omega_j\frac{\partial w}{\partial x_j},
\end{equation}
and
\begin{equation}\label{eq:3-85}
\frac{\partial}{\partial t}\nabla^2 w = g\alpha\!\left(\frac{\partial^2\Theta}{\partial x^2}+\frac{\partial^2\Theta}{\partial y^2}\right) + \nu\nabla^4 w - 2\Omega_j\frac{\partial\zeta}{\partial x_j},
\end{equation}
where $w$ and $\zeta$ are the $z$-components of the velocity and the vorticity, respectively.

Most of our discussions relating to equations~\eqref{eq:3-84} and~\eqref{eq:3-85} will be restricted to the case when $\boldsymbol{\Omega}$ and $\mathbf{g}$ are parallel; the case when they are not will be briefly considered in \S\,34.

When the axis of rotation coincides with the vertical, the relevant equations are
\begin{align}
\frac{\partial\Theta}{\partial t} &= \beta w + \kappa\nabla^2\Theta, \label{eq:3-86}\\[4pt]
\frac{\partial\zeta}{\partial t} &= 2\Omega\frac{\partial w}{\partial z} + \nu\nabla^2\zeta, \label{eq:3-87}
\end{align}
and
\begin{equation}\label{eq:3-88}
\frac{\partial}{\partial t}\nabla^2 w = g\alpha\!\left(\frac{\partial^2\Theta}{\partial x^2}+\frac{\partial^2\Theta}{\partial y^2}\right) + \nu\nabla^4 w - 2\Omega\frac{\partial\zeta}{\partial z},
\end{equation}
and we have to seek solutions of these equations which satisfy the boundary conditions described in Chapter~II, \S\,9\,(a).

\subsection*{(a) The analysis into normal modes}

Following the procedure in Chapter~II, \S\,10, we analyse the perturbations $w$, $\Theta$, and $\zeta$ into two-dimensional periodic waves; and considering disturbances characterized by a particular wave number $k$, we suppose that $w$, $\Theta$, and $\zeta$ have the forms assumed in equation~II~(90). Equations~\eqref{eq:3-86}--\eqref{eq:3-88} then become
\begin{align}
p\,\Theta &= \beta W + \kappa\!\left(\frac{d^2}{dz^2}-k^2\right)\!\Theta, \label{eq:3-89}\\[4pt]
p\,Z &= 2\Omega\frac{dW}{dz} + \nu\!\left(\frac{d^2}{dz^2}-k^2\right)\!Z, \label{eq:3-90}\\[4pt]
p\!\left(\frac{d^2}{dz^2}-k^2\right)\!W &= -g\alpha k^2\Theta + \nu\!\left(\frac{d^2}{dz^2}-k^2\right)^{\!2}\!W - 2\Omega\frac{dZ}{dz}. \label{eq:3-91}
\end{align}

Measuring lengths in the unit $d$ and letting
\begin{equation}\label{eq:3-92}
a = kd, \quad \sigma = pd^2/\nu, \quad \text{and} \quad \mathscr{P} = \nu/\kappa,
\end{equation}
we can reduce equations~\eqref{eq:3-89}--\eqref{eq:3-91} to the forms
\begin{align}
(D^2 - a^2 - \mathscr{P}\sigma)\,\Theta &= -\left(\frac{\beta d^2}{\kappa}\right)\!W, \label{eq:3-93}\\[4pt]
(D^2 - a^2 - \sigma)\,Z &= -\left(\frac{2\Omega}{\nu}\,d\right)\!DW, \label{eq:3-94}\\[4pt]
(D^2 - a^2)(D^2 - a^2 - \sigma)\,W - \left(\frac{2\Omega}{\nu}\,d^3\right)\!DZ &= \left(\frac{g\alpha d^2}{\nu}\right)\!a^2\Theta, \label{eq:3-95}
\end{align}
where $D = d/dz$ ($z$ being measured in the new unit). Solutions of equations~\eqref{eq:3-93}--\eqref{eq:3-95} must be sought which satisfy the boundary conditions given in equations~II~(95) and~(96).

\section{The case when instability sets in as stationary convection. A variational principle}\label{sec:3-26}

We shall find that in contrast to the simple B\'enard problem, the principle of the exchange of stabilities is not valid, generally, when rotation is present. However, it is not possible to state in simple analytical terms the necessary and sufficient conditions for its validity; it will appear that they are best stated in terms of the solution for the case of instability as stationary convection. For this reason, we shall first establish the conditions for this latter type of instability to occur before deciding whether instability will in fact occur this way.

When instability sets in as ordinary convection, the marginal state will be characterized by $\sigma = 0$, and the basic equations reduce to (cf.\ equations~\eqref{eq:3-93}--\eqref{eq:3-95})
\begin{align}
(D^2 - a^2)\,\Theta &= -\left(\frac{\beta d^2}{\kappa}\right)\!W, \label{eq:3-96}\\[4pt]
(D^2 - a^2)\,Z &= -\left(\frac{2\Omega}{\nu}\,d\right)\!DW, \label{eq:3-97}
\end{align}
and
\begin{equation}\label{eq:3-98}
(D^2 - a^2)^2 W - \left(\frac{2\Omega}{\nu}\,d^3\right)\!DZ = \left(\frac{g\alpha d^2}{\nu}\right)\!a^2\Theta.
\end{equation}

We can eliminate $Z$ and $\Theta$ from the last equation by operating $(D^2 - a^2)$ on it. We find:
\begin{equation}\label{eq:3-99}
(D^2 - a^2)^3 W + \mathscr{T}\,D^2 W = -Ra^2 W,
\end{equation}
where $R$ is the Rayleigh number as usually defined, and
\begin{equation}\label{eq:3-100}
\mathscr{T} = \frac{4\Omega^2}{\nu^2}\,d^4
\end{equation}
is the Taylor number.

The boundary conditions with respect to which we must seek solutions of the foregoing equations are (cf.\ equations~II~(95) and~(96))
\begin{equation}\label{eq:3-101}
W = 0 \quad \text{and} \quad \Theta = 0 \quad \text{for } z = 0 \text{ and } 1,
\end{equation}
\begin{equation}\label{eq:3-102}
\text{either}\quad Z = 0 \quad \text{and} \quad DW = 0 \quad \text{(on a rigid surface)},
\end{equation}
\begin{equation*}
\text{or}\quad DZ = 0 \quad \text{and} \quad D^2 W = 0 \quad \text{(on a free surface)}.
\end{equation*}

In view of equation~\eqref{eq:3-98} we can replace the conditions~\eqref{eq:3-101} by
\begin{equation}\label{eq:3-103}
W = 0 \quad \text{and} \quad (D^2 - a^2)^2 W - \left(\frac{2\Omega}{\nu}\,d^3\right)\!DZ = 0 \quad \text{for } z = 0 \text{ and } 1.
\end{equation}

Since the boundary conditions~\eqref{eq:3-103} involve $Z$, it is clear that we cannot treat equation~\eqref{eq:3-99} independently of equation~\eqref{eq:3-97} for $Z$; the order of the system we are effectively dealing with is, therefore, eight not six.

Equations~\eqref{eq:3-97} and~\eqref{eq:3-99}, together with the boundary conditions~\eqref{eq:3-102} and~\eqref{eq:3-103}, clearly constitute a characteristic value problem for $R$ for given $a^2$ and $\mathscr{T}$; and the problem of determining the critical Rayleigh number for the onset of instability as stationary convection reduces to the following.

We must first determine the lowest characteristic value of $R$, through the solution of equations~\eqref{eq:3-97}, \eqref{eq:3-99}, \eqref{eq:3-102}, and~\eqref{eq:3-103}, as a function of $a^2$ for a given $\mathscr{T}$ and then find the minimum of this function; and the minimum so determined is the critical Rayleigh number for the onset of stationary convection for the given~$\mathscr{T}$.

\subsection*{(a) A variational principle}

As in the case of the simple B\'enard problem, we can formulate the present characteristic value problem also in terms of a variational principle.

Letting
\begin{equation}\label{eq:3-104}
F = (D^2 - a^2)^2 W - \left(\frac{2\Omega}{\nu}\,d^3\right)\!DZ,
\end{equation}
we can rewrite the differential equations governing $W$ and $Z$ in the forms
\begin{equation}\label{eq:3-105}
(D^2 - a^2)F = -Ra^2 W
\end{equation}
and
\begin{equation}\label{eq:3-106}
(D^2 - a^2)Z = -\left(\frac{2\Omega}{\nu}\,d\right)\!DW.
\end{equation}

The boundary conditions~\eqref{eq:3-103} require that
\begin{equation}\label{eq:3-107}
W = F = 0 \quad \text{for } z = 0 \text{ and } 1.
\end{equation}

Now multiply equation~\eqref{eq:3-105} by $F$ and integrate over the range of $z$. After an integration by parts, the left-hand side gives
\begin{equation}\label{eq:3-108}
\int_0^1 F(D^2 - a^2)F\,dz = -\int_0^1 [(DF)^2 + a^2 F^2]\,dz,
\end{equation}
the integrated part vanishing on account of the boundary condition on $F$; and the right-hand side of the equation requires us to consider
\begin{equation}\label{eq:3-109}
\int_0^1 W F\,dz = \int_0^1 W(D^2 - a^2)^2 W\,dz - \frac{2\Omega}{\nu}\,d^3 \int_0^1 W\,DZ\,dz.
\end{equation}

After two integrations by parts, the first of the two integrals on the right-hand side of~\eqref{eq:3-109} becomes (cf.\ equation~II~(120))
\begin{equation}\label{eq:3-110}
\int_0^1 W(D^2 - a^2)^2 W\,dz = \int_0^1 [(D^2 - a^2)W]^2\,dz,
\end{equation}
while the second, after an integration by parts, gives (cf.\ equations~\eqref{eq:3-103} and~\eqref{eq:3-106})
\begin{equation}\label{eq:3-111}
-\frac{2\Omega}{\nu}\,d^3\int_0^1 W\,DZ\,dz = \frac{2\Omega}{\nu}\,d^3\int_0^1 Z\,DW\,dz = -d^2\int_0^1 Z(D^2-a^2)Z\,dz.
\end{equation}

Remembering that on a bounding surface either $Z$ or $DZ$ vanishes, we obtain after a further integration by parts
\begin{equation}\label{eq:3-112}
-\frac{2\Omega}{\nu}\,d^3\int_0^1 W\,DZ\,dz = d^2 \int_0^1 [(DZ)^2 + a^2 Z^2]\,dz.
\end{equation}

Combining equations~\eqref{eq:3-109}, \eqref{eq:3-110}, and~\eqref{eq:3-112}, we have
\begin{equation}\label{eq:3-113}
\int_0^1 W F\,dz = \int_0^1 [(D^2 - a^2)W]^2\,dz + d^2 \int_0^1 [(DZ)^2 + a^2 Z^2]\,dz.
\end{equation}

The result of multiplying equation~\eqref{eq:3-105} by $F$ and integrating over $z$ is
\begin{equation}\label{eq:3-114}
R = \frac{\displaystyle\int_0^1 [(DF)^2 + a^2 F^2]\,dz}{a^2\displaystyle\int_0^1\bigl\{[(D^2 - a^2)W]^2 + d^2[(DZ)^2 + a^2 Z^2]\bigr\}\,dz} = \frac{I_1}{a^2 I_2} \quad \text{(say)}.
\end{equation}

This formula expresses $R$ as a ratio of two positive definite integrals. Consider now the effect on $R$ of variations $\delta W$ and $\delta Z$ in $W$ and $Z$ compatible with the boundary conditions on $W$ and $Z$. To the first order, we have
\begin{equation}\label{eq:3-115}
\delta R = \frac{1}{a^2 I_2}\!\left(\delta I_1 - \frac{I_1}{I_2}\,\delta I_2\right) = \frac{1}{a^2 I_2}(\delta I_1 - Ra^2\,\delta I_2).
\end{equation}

$\delta I_1$ and $\delta I_2$ are the corresponding variations in $I_1$ and $I_2$:
\begin{align}
\delta I_1 &= 2\int_0^1 [(DF)(D\,\delta F) + a^2 F\,\delta F]\,dz, \notag\\[6pt]
\delta I_2 &= 2\int_0^1 \bigl\{(D^2 - a^2)W [(D^2 - a^2)\,\delta W] + \notag\\
&\qquad\qquad + d^2[(DZ)(D\,\delta Z) + a^2 Z\,\delta Z]\bigr\}\,dz. \label{eq:3-116}
\end{align}

Making use of the boundary conditions which $F$, $\delta F$, $W$, $\delta W$, $Z$, and $\delta Z$ satisfy, we can reduce the expressions for $\delta I_1$ and $\delta I_2$ by one or more integrations by parts. Thus
\begin{equation}\label{eq:3-117}
\delta I_1 = -2\int_0^1 \delta F(D^2 - a^2)F\,dz,
\end{equation}
and
\begin{align}
\delta I_2 &= 2\int_0^1 W(D^2 - a^2)^2\,\delta W - 2d^2\int_0^1 \delta Z(D^2 - a^2)Z\,dz \notag\\
&= 2\int_0^1 W[(D^2 - a^2)^2\,\delta W + \frac{4\Omega^2}{\nu^2}d^4\int_0^1 \delta Z\,DW\,dz \notag\\
&= 2\int_0^1 W\!\left[(D^2 - a^2)^2\,\delta W - \left(\frac{2\Omega}{\nu}\,d^3\right)\!D\,\delta Z\right]\!dz \label{eq:3-118}\\
&= 2\int_0^1 W\,\delta F\,dz. \notag
\end{align}

Now combining equations~\eqref{eq:3-115}, \eqref{eq:3-117}, and~\eqref{eq:3-118}, we have
\begin{equation}\label{eq:3-119}
\delta R = -\frac{2}{a^2 I_2}\int_0^1 \delta F[(D^2 - a^2)F + Ra^2 W]\,dz,
\end{equation}
where it may be recalled that equation~\eqref{eq:3-104} defining $F$ and equation~\eqref{eq:3-106} relating $Z$ and $W$ have been explicitly used in the reductions.

From equation~\eqref{eq:3-119} it follows that
\begin{equation}\label{eq:3-120}
\delta R = 0 \quad \text{if}\quad (D^2 - a^2)F = -Ra^2 W;
\end{equation}
\emph{and conversely if $\delta R = 0$ for any arbitrary variation,}
\[
\delta F = (D^2 - a^2)^2\,\delta W - (2\Omega d^3/\nu)\,D\,\delta Z,
\]
\emph{compatible with the boundary conditions of the problem, then equation~\eqref{eq:3-120} must hold and the function $W$ in terms of which $R$ was initially calculated must have been a solution of the characteristic value problem.}

The foregoing arguments establish the stationary property of the characteristic values when they are regarded as given by equation~\eqref{eq:3-114}. And it can be shown, by a procedure exactly analogous to that used in connexion with the first variational principle for the simple B\'enard problem (Chapter~II, \S\,13\,(a)), that the lowest characteristic value of $R$ is, indeed, a true minimum of the functional~\eqref{eq:3-114}. We are thus enabled to formulate the following variational procedure for solving equations~\eqref{eq:3-97} and~\eqref{eq:3-99} (for any assigned $a^2$ and $\mathscr{T}$) and satisfying the boundary conditions of the problem.

Assume for $F$ an expansion involving one or more parameters $A_p$ which vanishes for $z = 0$ and~1. With the chosen form of $F$, determine $W$ and $Z$ as solutions of the equations
\begin{equation}\label{eq:3-121}
(D^2 - a^2)^2 W - \left(\frac{2\Omega}{\nu}\,d^3\right)\!DZ = F
\end{equation}
and
\begin{equation}\label{eq:3-122}
(D^2 - a^2)Z = -\left(\frac{2\Omega}{\nu}\,d\right)\!DW,
\end{equation}
which satisfy a total of six boundary conditions on $W$ and $Z$ at $z = 0$ and~1. Since equations~\eqref{eq:3-121} and~\eqref{eq:3-122} are together of order six, there will be just enough constants of integration in the general solution to satisfy all the boundary conditions. Having solved for $W$ and $Z$ in this manner, evaluate $R$ according to the formula~\eqref{eq:3-114} and minimize it with respect to the parameters $A_p$. In this way, we shall obtain the `best' value of $R$ for the chosen form of $F$. We shall see in the following section that even with the simplest trial function for $F$, we can reach quite high precision in the deduced values of $R$.

Finally, it may be noted that according to equations~\eqref{eq:3-121} and~\eqref{eq:3-122}, the equation for $W$ that must be solved when using the variational method is
\begin{equation}\label{eq:3-123}
(D^2 - a^2)^3 W + \mathscr{T}\,D^2 W = (D^2 - a^2)F,
\end{equation}
where $\mathscr{T}$ is the Taylor number.

\section{Solutions for the case when instability sets in as stationary convection}\label{sec:3-27}

We shall obtain the solution for the following three cases: when both the bounding surfaces are free; when both the bounding surfaces are rigid; and when one bounding surface is rigid and the other is free.

\subsection*{(a) The solution for two free boundaries}

In this case the boundary conditions~\eqref{eq:3-102} and~\eqref{eq:3-103} require
\begin{equation}\label{eq:3-124}
W = D^2 W = D^4 W = 0 \quad \text{and} \quad DZ = 0 \quad \text{for } z = 0 \text{ and } 1.
\end{equation}

From the equation satisfied by $W$ (namely~\eqref{eq:3-99}), it follows that $D^6 W = 0$ for $z = 0$ and~1. By differentiating equation~\eqref{eq:3-99} even numbers of times, we can successively conclude that all the even derivatives of $W$ must, therefore, be
\begin{equation}\label{eq:3-125}
W = A\sin n\pi z,
\end{equation}
where $A$ is a constant and $n$ is an integer. The corresponding solution for $Z$ is
\begin{equation}\label{eq:3-126}
Z = A\!\left(\frac{2\Omega}{\nu}\,d\right)\!\frac{n\pi}{n^2\pi^2 + a^2}\cos n\pi z.
\end{equation}

Substitution of the solution~\eqref{eq:3-125} in equation~\eqref{eq:3-99} leads to the characteristic equation
\begin{equation}\label{eq:3-127}
R = \frac{1}{a^2}[(n^2\pi^2 + a^2)^3 + n^2\pi^2\,\mathscr{T}].
\end{equation}

For a given $a^2$, the lowest characteristic value for $R$ occurs when $n = 1$; then,
\begin{equation}\label{eq:3-128}
R = \frac{1}{a^2}[(\pi^2 + a^2)^3 + \pi^2\,\mathscr{T}];
\end{equation}
and this equation has to be interpreted in the same way as equation~II~(193) was interpreted in Chapter~II, \S\,15\,(a).

Letting
\begin{equation}\label{eq:3-129}
a^2 = \pi^2 x,
\end{equation}
we can rewrite equation~\eqref{eq:3-128} in the form
\begin{equation}\label{eq:3-130}
R = \pi^4\!\left[(1+x)^3 + \frac{\mathscr{T}}{\pi^4}\right]\!\frac{1}{x}.
\end{equation}

\begin{table}[ht]
\centering
\caption{Critical Rayleigh numbers and wave numbers of the unstable modes at marginal stability for the onset of stationary convection when both bounding surfaces are free}
\label{tab:VII}
\begin{tabular}{ccccccc}
\hline
$\mathscr{T}$ & $a_c$ & $R_c$ & $\mathscr{T}$ & $a_c$ & $R_c$ \\
\hline
0 & $2{\cdot}233$ & $6{\cdot}575\times 10^2$ & $3\times 10^4$ & $10{\cdot}45$ & $4{\cdot}237\times 10^4$ \\
$10$ & $2{\cdot}270$ & $6{\cdot}771\times 10^2$ & $10^5$ & $12{\cdot}85$ & $9{\cdot}232\times 10^4$ \\
$10^2$ & $2{\cdot}594$ & $8{\cdot}263\times 10^2$ & $10^6$ & $19{\cdot}02$ & $4{\cdot}147\times 10^5$ \\
$5\times 10^2$ & $3{\cdot}228$ & $1{\cdot}473\times 10^3$ & $10^7$ & $28{\cdot}02$ & $1{\cdot}897\times 10^6$ \\
$10^3$ & $3{\cdot}710$ & & $10^8$ & $41{\cdot}20$ & $8{\cdot}536\times 10^6$ \\
$10^4$ & $4{\cdot}211$ & $2{\cdot}009\times 10^4$ & $10^9$ & $60{\cdot}53$ & $4{\cdot}017\times 10^7$ \\
$5\times 10^3$ & $5{\cdot}021$ & $5{\cdot}999\times 10^3$ & $10^{10}$ & $88{\cdot}87$ & $1{\cdot}879\times 10^8$ \\
$5\times 10^3$ & $5{\cdot}665$ & $3{\cdot}377\times 10^4$ & $10^{11}$ & $130{\cdot}46$ & $8{\cdot}702\times 10^8$ \\
$3\times 10^4$ & $6{\cdot}054$ & $1{\cdot}013\times 10^4$ & $10^{12}$ & $191{\cdot}31$ & $4{\cdot}037\times 10^9$ \\
$10^4$ & $8{\cdot}245$ & & & & \\
\hline
\end{tabular}
\end{table}

As a function of $x$, $R$ given by equation~\eqref{eq:3-130} attains its minimum when
\begin{equation}\label{eq:3-131}
2x^3 + 3x^2 = 1 + \mathscr{T}/\pi^4.
\end{equation}

With $x$ determined as the root of this cubic equation, equation~\eqref{eq:3-130} will give the required critical Rayleigh number $R_c$. Values of $R_c$ determined in this fashion for various values of $\mathscr{T}$ are given in Table~VII; the wave numbers characterizing the marginal states are also given. The results are further illustrated in Figs.~21 and~22. The inhibiting effect of rotation on the onset of instability is apparent from these results.

For $\mathscr{T}/\pi^4$ sufficiently large, the required root of equation~\eqref{eq:3-131} tends to
\begin{equation}\label{eq:3-132}
x_{\min} \to \left(\frac{\mathscr{T}}{2\pi^4}\right)^{\!1/3} \quad (\mathscr{T}\to\infty).
\end{equation}

The corresponding asymptotic behaviours of $R_c$ and $a_{\min}$ are:
\begin{equation}\label{eq:3-133}
R_c \to 3\pi^4\!\left(\frac{\mathscr{T}}{2\pi^4}\right)^{\!2/3} = 8{\cdot}6956\,\mathscr{T}^{2/3} \quad (\mathscr{T}\to\infty).
\end{equation}
and
\[
a_{\min} \to (\tfrac{1}{2}\pi^2\,\mathscr{T})^{1/6} = 1{\cdot}3048\,\mathscr{T}^{1/6}.
\]

\figurePlaceholder{21}{The $(R_c,\mathscr{T})$-relations for the three cases (i) both bounding surfaces rigid, (ii) one bounding surface rigid and the other free, and (iii) both bounding surfaces free; the curves labelled $a$, $b$, and $c$ are the relations for the onset of ordinary cellular convection for the three cases, respectively. The curves labelled $a'A$, $b'B$, and $c'C'$ are the corresponding relations for the onset of overstability for $\mathscr{P} = 0{\cdot}025$. At $c$ (respectively $c'$) we have a change from one type of instability to another as $\mathscr{T}$ increases.}

\figurePlaceholder{22}{The $(a_c,\mathscr{T})$-relations for the three cases (i) both bounding surfaces rigid, (ii) one bounding surface rigid and the other free, and (iii) both bounding surfaces free; the curves labelled $a$, $b$, and $c$ are the relations for the onset of ordinary cellular convection for the three cases, respectively. The curves labelled $a'A$, $B$, and $C'$ are the corresponding relations for the onset of overstability for $\mathscr{P} = 0{\cdot}025$.}

Substituting for $R$ and $\mathscr{T}$ in accordance with their definitions, we find that the formula which determines the critical temperature gradient for the onset of stationary convection, when $\mathscr{T}\to\infty$, is
\begin{equation}\label{eq:3-134}
g\alpha\beta_c \to 21{\cdot}911\!\left(\frac{\Omega}{d}\right)^{\!4/3}\!\kappa^{-1} \quad (\Omega\to\infty;\ \text{or}\ \nu\to 0);
\end{equation}
this should be contrasted with the formula,
\begin{equation}\label{eq:3-135}
g\alpha\beta_c = \text{constant}\ \kappa d^{-4},
\end{equation}
valid in the absence of rotation. The dependence of $\beta_c$ on an inverse power of $\nu$, for $\Omega$ fixed and $\nu\to 0$, implies that an \emph{inviscid, ideal fluid, in rotation, is stable with respect to the onset of stationary convection for all adverse temperature gradients}. This is clearly a consequence of the Taylor--Proudman theorem.

\subsection*{(b) The solution for two rigid boundaries}

We shall obtain the solution for this case by the variational method. In view of the symmetry of this problem with respect to the bounding planes, we shall find it convenient to translate the origin of $z$ to be midway between the two planes. The fluid will then be confined between $z = \pm\frac{1}{2}$ and we shall have to seek solutions of equations~\eqref{eq:3-120}--\eqref{eq:3-123} which satisfy the boundary conditions
\begin{equation}\label{eq:3-136}
F = W = DW = Z = 0 \quad \text{for } z = \pm\tfrac{1}{2}.
\end{equation}

It is apparent from the equations and the boundary conditions that the proper solutions of the problem fall into two non-combining groups; these consist of solutions which are even in $W$ and odd in $Z$ and solutions which are odd in $W$ and even in $Z$; we shall call these the even and the odd solutions, respectively. It is also clear that the lowest characteristic value of $R$ will occur among the even solutions; we shall accordingly consider such solutions.

Since $F$ is assumed to be even and is required to vanish at $z = \pm\frac{1}{2}$, we can expand it in a cosine series in the form
\begin{equation}\label{eq:3-137}
F = \sum_m A_m\cos[(2m+1)\pi z],
\end{equation}
where the summation over $m$ may be assumed to run from zero to infinity; but this is not necessary since we shall consider the coefficients $A_m$ as variational parameters.

With the chosen form for $F$, the equation to be solved for $W$ is (cf.\ equation~\eqref{eq:3-123})
\begin{equation}\label{eq:3-138}
[(D^2 - a^2)^3 + \mathscr{T}\,D^2]\,W = -\sum_m A_m\, c_{m+1}\cos[(2m+1)\pi z],
\end{equation}
where
\begin{equation}\label{eq:3-139}
c_{m+1} = (2m+1)^2\pi^2 + a^2.
\end{equation}

In view of the linearity of equation~\eqref{eq:3-138}, we can express $W$ and $Z$ as sums of the form
\begin{equation}\label{eq:3-140}
W = \sum_m A_m\, W_m \quad \text{and} \quad Z = \sum_m A_m\, Z_m,
\end{equation}
where $W_m$ and $Z_m$ are solutions of the equations
\begin{equation}\label{eq:3-141}
(D^2 - a^2)^3 W_m + \mathscr{T}\,D^2 W_m = -c_{m+1}\cos[(2m+1)\pi z]
\end{equation}
and
\begin{equation}\label{eq:3-142}
(D^2 - a^2)Z_m = -\left(\frac{2\Omega}{\nu}\,d\right)\!DW_m.
\end{equation}

The variational principle formulated in \S\,26\,(a) is equivalent to minimising
\begin{equation}\label{eq:3-143}
\int_{-1/2}^{+1/2} [(DF)^2 + a^2 F^2]\,dz
\end{equation}
for such variations of the parameters (the $A_m$'s in the present connexion) which preserve the constancy of
\begin{equation}\label{eq:3-144}
\int_{-1/2}^{+1/2} WF\,dz.
\end{equation}

By arguments similar to those used in Chapter~II, \S\,17, it can be shown, by the use of a Lagrangian undetermined multiplier, that the application of this principle is exactly equivalent to solving a secular determinant which will arise from a Fourier analysis of the equation
\begin{equation}\label{eq:3-145}
(D^2 - a^2)\sum_m A_m\cos[(2m+1)\pi z] = -Ra^2\sum_m A_m\, W_m.
\end{equation}

Returning to equation~\eqref{eq:3-141}, we observe that the general solution of this equation, which is even, can be expressed in the form
\begin{equation}\label{eq:3-146}
W_m = c_{m+1}\,\gamma_{2m+1}\cos[(2m+1)\pi z] + \sum_{j=1}^{3} B_j^{(m)}\cosh q_j\, z,
\end{equation}
where
\begin{equation}\label{eq:3-147}
\frac{1}{\gamma_{2m+1}} = [(2m+1)^2\pi^2 + a^2]^3 + (2m+1)^2\pi^2\,\mathscr{T} = c_{2m+1}^3 + (2m+1)^2\pi^2\,\mathscr{T},
\end{equation}
the $B_j^{(m)}$'s ($j = 1,2,3$) are constants of integration, and the $q_j$'s ($j = 1,2,3$) are the three roots of the cubic equation
\begin{equation}\label{eq:3-148}
(q^2 - a^2)^3 + \mathscr{T}\,q^2 = 0.
\end{equation}

An identity which readily follows from equation~\eqref{eq:3-146} is
\begin{equation}\label{eq:3-149}
\frac{1}{\gamma_{2m+1}} = \frac{1}{\Delta}[(2m+1)^2\pi^2 + q_j^2];
\end{equation}
for,
\begin{equation}\label{eq:3-150}
(q^2 - a^2)^3 + \mathscr{T}\,q^2 = \prod_{j=1}^{3}(q^2 - q_j^2)
\end{equation}
is, in reality, an identity.

The solution for $Z_m$ corresponding to the solution~\eqref{eq:3-146} for $W_m$ is
\begin{align}
Z_m &= -\left(\frac{2\Omega}{\nu}\,d\right)\!\bigl\{(2m+1)\pi\gamma_{2m+1}\sin[(2m+1)\pi z] + \notag\\
&\qquad\qquad + \sum_{j=1}^{3} B_j^{(m)}\frac{q_j}{x_j}\sinh q_j\, z\bigr\}, \label{eq:3-151}
\end{align}
where
\begin{equation}\label{eq:3-152}
x_j = q_j^2 - a^2.
\end{equation}

The boundary conditions $W = DW = Z = 0$ for $z = \pm\frac{1}{2}$ require
\begin{align}
&\sum_{j=1}^{3} B_j^{(m)}\cosh\tfrac{1}{2}q_j = 0, \notag\\[4pt]
&\sum_{j=1}^{3} B_j^{(m)}\,q_j\sinh\tfrac{1}{2}q_j = (-1)^m(2m+1)\pi\gamma_{2m+1}\, c_{m+1}, \label{eq:3-153}\\[4pt]
&\sum_{j=1}^{3} B_j^{(m)}\frac{q_j}{x_j}\sinh\tfrac{1}{2}q_j = (-1)^{m+1}(2m+1)\pi\gamma_{2m+1}. \notag
\end{align}

On solving these equations, we find that
\begin{align}
B_1^{(m)} &= (-1)^m(2m+1)\pi\gamma_{2m+1}\bigl\{q_1\bigl(1+\tfrac{c_{m+1}}{x_1}\bigr)\coth\tfrac{1}{2}q_1\, s - \notag\\
&\qquad\qquad -q_2\bigl(1+\tfrac{c_{m+1}}{x_2}\bigr)\coth\tfrac{1}{2}q_2\,\Delta\operatorname{cosech}\tfrac{1}{2}q_1\bigr\}, \label{eq:3-154}
\end{align}
where
\begin{equation}\label{eq:3-155}
\frac{1}{\Delta} = \frac{q_1 q_2}{x_1 x_2}(x_2 - x_1)\coth\tfrac{1}{2}q_1 + \frac{q_1 q_2}{x_1 x_2}(x_1 - x_2)\coth\tfrac{1}{2}q_2 + \frac{q_1 q_2}{x_1 x_2}(x_2 - x_1)\coth\tfrac{1}{2}q_1,
\end{equation}
and $B_2^{(m)}$ and $B_3^{(m)}$ are given by similar expressions which can be obtained by cyclically permuting the $q_j$'s and the $x_j$'s in the solution~\eqref{eq:3-154}.

The characteristic equation is now obtained from (cf.\ equations~\eqref{eq:3-145} and~\eqref{eq:3-146})
\begin{equation}\label{eq:3-156}
\sum_m A_m\, c_{m+1}\cos[(2m+1)\pi z] = Ra^2\sum_m A_m\!\left[c_{m+1}\,\gamma_{2m+1}\cos[(2m+1)\pi z] + \sum_{j=1}^{3} B_j^{(m)}\cosh q_j\, z\right].
\end{equation}

Multiplying this equation by $\cos[(2n+1)\pi z]$ and integrating over the range of $z$, we obtain
\begin{equation}\label{eq:3-157}
\tfrac{1}{2}c_{m+1}\, A_m = Ra^2\!\left[\tfrac{1}{2}c_{m+1}\,\gamma_{2m+1}\, A_m + \sum_n (n|m)\,A_n\right] \quad (n = 0,1,2,\ldots),
\end{equation}
where
\begin{equation}\label{eq:3-158}
(n|m) = \sum_{j=1}^{3} B_j^{(m)} \int_{-1/2}^{+1/2} \cosh q_j\, z\cos[(2n+1)\pi z]\,dz
= 4m\pi(-1)^n\sum_{j=1}^{3} \frac{B_j^{(m)}\cosh\tfrac{1}{2}q_j}{(2n+1)^2\pi^2 + q_j^2}.
\end{equation}

Equation~\eqref{eq:3-158} provides a set of linear homogeneous equations for the constants $A_m$; this same set of equations ensures that expression~\eqref{eq:3-143} is a minimum for all variations of the $A_m$'s which preserve the constancy of~\eqref{eq:3-144}. In this latter formulation $Ra^2$ is the Lagrangian undetermined multiplier.

The determinant of the system of equations represented by~\eqref{eq:3-157} must vanish; and this provides the characteristic equation for $R$. Thus,
\begin{equation}\label{eq:3-159}
\left\|\tfrac{1}{2}c_{m+1}\!\left[\frac{1}{Ra^2} - \gamma_{2m+1}\right]\!\delta_{mn} - (n|m)\right\| = 0.
\end{equation}

Substituting for the $B_j^{(m)}$'s in accordance with equation~\eqref{eq:3-154} and making use of the identity~\eqref{eq:3-149}, we find after some minor reductions that the expression for $(n|m)$ given in~\eqref{eq:3-158} becomes
\begin{align}
(n|m) &= 2(-1)^{m+n}(2n+1)(2m+1)\pi^2\gamma_{2m+1}\,\gamma_{2n+1}\,\Delta\prod_{j=1}^{3}\coth\tfrac{1}{2}q_j\times \notag\\
&\quad\times\Bigl\{q_1(q_2^2 - q_3^2)\bigl[(2n+1)^2\pi^2 + q_1^2\bigr]\bigl(1+\tfrac{c_{m+1}}{x_1}\bigr)\tanh\tfrac{1}{2}q_1 + \notag\\
&\qquad + q_2(q_1^2 - q_3^2)(2n+1)^2\pi^2 + q_2^2\bigr]\bigl(1+\tfrac{c_{m+1}}{x_2}\bigr)\tanh\tfrac{1}{2}q_2 + \notag\\
&\qquad + q_3(q_1^2 - q_2^2)\bigl[(2n+1)^2\pi^2 + q_3^2\bigr]\bigl(1+\tfrac{c_{m+1}}{x_3}\bigr)\tanh\tfrac{1}{2}q_3\Bigr\}. \label{eq:3-160}
\end{align}

Recalling the definitions of $c_{m+1}$ and $x_j$ (equations~\eqref{eq:3-139} and~\eqref{eq:3-152}), we can rewrite the foregoing in the form
\begin{align}
(n|m) &= 2(-1)^{m+n}(2n+1)(2m+1)\pi^2\gamma_{2m+1}\,\gamma_{2n+1}\,\Delta\prod_{j=1}^{3}\coth\tfrac{1}{2}q_j\times \notag\\
&\quad\times\Bigl\{(x_2 - x_3)(c_{m+1} + x_1)\bigl(1+\tfrac{c_{m+1}}{x_1}\bigr)\,q_1\tanh\tfrac{1}{2}q_1 + \notag\\
&\qquad + (x_1 - x_3)(c_{m+1} + x_2)\bigl(1+\tfrac{c_{m+1}}{x_2}\bigr)\,q_2\tanh\tfrac{1}{2}q_2 + \notag\\
&\qquad + (x_1 - x_2)(c_{m+1} + x_3)\bigl(1+\tfrac{c_{m+1}}{x_3}\bigr)\,q_3\tanh\tfrac{1}{2}q_3\Bigr\}. \label{eq:3-161}
\end{align}

This expression for the matrix element $(n|m)$ is manifestly symmetric in $n$ and $m$: a fact which reflects the basic self-adjoint character of the underlying problem.

The expression for $(n|m)$ can be further reduced. First, we observe that according to equations~\eqref{eq:3-148} and~\eqref{eq:3-152} the $x_j$'s are the roots of the equation
\begin{equation}\label{eq:3-162}
x^3 + \mathscr{T}\,x + \mathscr{T}\,a^2 = 0.
\end{equation}

This equation allows one real root and a pair of complex conjugate roots; let these be
\begin{equation}\label{eq:3-163}
x \quad \text{and} \quad X \pm iY.
\end{equation}

Denote the corresponding roots of $q$ by
\begin{equation}\label{eq:3-164}
q = \sqrt{a^2 + x} \quad \text{and} \quad \alpha_1 \pm i\alpha_2 = \sqrt{a^2 + X \pm iY}.
\end{equation}

On substituting for the $x_j$'s and the $q_j$'s in accordance with these definitions, we find that the expression~\eqref{eq:3-161} for $(n|m)$ can be reduced to the form
\begin{align}
(n|m) &= 2(-1)^{m+1}(2n+1)(2m+1)\pi^2\gamma_{2m+1}\,\gamma_{2n+1}\,\Phi^{(0)}\times \notag\\
&\quad\times\Bigl\{Y(c_{2m+1}+x)\bigl(1+\tfrac{c_{2m+1}}{x}\bigr)\tanh\tfrac{1}{2}q + \notag\\
&\qquad + \tfrac{c_{2n+1}\,c_{2m+1}}{X^2+Y^2}\,\Phi_1^{(c)} - (c_{2n+1}+c_{2m+1})\,\Phi_2^{(c)} + \Phi_3^{(c)}\Bigr\}, \label{eq:3-165}
\end{align}
where $\Phi^{(0)}, \ldots, \Phi_3^{(c)}$ are functions of $q$, $\alpha_1$, $\alpha_2$, $x$, $X$, and $Y$ defined as follows.

Let
\begin{align}
\phi_1^{(c)} &= \frac{\alpha_1\sinh\alpha_1 - \alpha_2\sin\alpha_2}{\cosh\alpha_1 + \cos\alpha_2}, \qquad
\psi_1^{(c)} = \frac{\alpha_2\sinh\alpha_1 + \alpha_1\sin\alpha_2}{\cosh\alpha_1 + \cos\alpha_2}, \notag\\[6pt]
\phi_2^{(c)} &= \frac{\alpha_1\sinh\alpha_1 + \alpha_2\sin\alpha_2}{\cosh\alpha_1 - \cos\alpha_2}, \qquad
\psi_2^{(c)} = \frac{\alpha_2\sinh\alpha_1 - \alpha_1\sin\alpha_2}{\cosh\alpha_1 - \cos\alpha_2}, \label{eq:3-166}
\end{align}
then
\begin{align}
\Phi^{(c)} &= \coth\tfrac{1}{2}q\,(\sinh^2\!\tfrac{1}{2}\alpha_1 + \sin^2\!\tfrac{1}{2}\alpha_2)^{-1}\times \notag\\
&\quad\times\left\{q\frac{Y\phi_1^{(c)} - X\psi_1^{(c)}}{X^2+Y^2} + \frac{\phi_1^{(c)}}{z}\right\} - \frac{\alpha_1+\alpha_2}{X^2+Y^2}\,Y\coth\tfrac{1}{2}q, \label{eq:3-167}
\end{align}
\begin{align}
\Phi_1^{(c)} &= (X^2 - Y^2 - xX)\psi_1^{(c)} - Y(2X-x)\phi_1^{(c)}, \notag\\[4pt]
\Phi_2^{(c)} &= Y\phi_1^{(c)} - (X-x)\psi_1^{(c)}, \notag\\[4pt]
\Phi_3^{(c)} &= (X^2 + Y^2 - xX)\psi_1^{(c)} - xY\phi_1^{(c)}. \label{eq:3-168}
\end{align}

By solving the determinantal equation~\eqref{eq:3-159} by including successively more rows and columns, we can evaluate the required characteristic values of $R$ with increasing precision. The results of such calculations are summarized in Table~VIII and illustrated in Figs.~21 and~22 (see p.~96). It will be seen that in no case is it really necessary to go higher than the second approximation.

\begin{table}[ht]
\centering
\caption{Critical Rayleigh numbers and related constants for the case when both bounding surfaces are rigid and the onset of instability is as stationary convection}
\label{tab:VIII}
\begin{tabular}{cccccccc}
\hline
& & \multicolumn{3}{c}{$R_c$} & \multicolumn{1}{c}{Second} & \multicolumn{2}{c}{Third approximation} \\
\cline{3-5}\cline{7-8}
$\mathscr{T}$ & $a_c$ & First & Second & Third & approxi- & & \\
& & approx. & approx. & approx. & mation $A_1/A_0$ & $A_1/A_0$ & $A_2/A_1$ \\
\hline
$10$ & $3{\cdot}10$ & $1{\cdot}720\times10^3$ & & $+0{\cdot}00881$ & & \\
$100$ & $3{\cdot}15$ & $1{\cdot}764\times10^3$ & & $+0{\cdot}00845$ & & \\
$500$ & $3{\cdot}45$ & $1{\cdot}948\times10^3$ & $1{\cdot}9405\times10^3$ & $+0{\cdot}00840$ & & \\
$1{,}000$ & $3{\cdot}65$ & $1{\cdot}945\times10^3$ & & $+0{\cdot}02349$ & $+0{\cdot}00091$ \\
$3{,}000$ & $3{\cdot}95$ & & & & \\
$5{,}000$ & $4{\cdot}80$ & & $3{\cdot}4688\times10^4$ & $+0{\cdot}00233$ & & \\
$10{,}000$ & & & & & & \\
$10^5$ & & & & $-0{\cdot}00403$ & & \\
$10^6$ & $10{\cdot}65$ & $1{\cdot}737\times10^5$ & & $-0{\cdot}00493$ & & \\
$10^7$ & $14{\cdot}5$ & $3{\cdot}513\times10^5$ & & $+0{\cdot}00055$ & & \\
$10^8$ & $44{\cdot}5$ & $3{\cdot}4698\times10^6$ & $3{\cdot}4574\times10^6$ & $-0{\cdot}00731$ & $-0{\cdot}07798$ & $+0{\cdot}04193$ \\
\hline
\end{tabular}
\end{table}

\subsection*{(c) The solution for one rigid and one free boundary}

In this case the conditions to be satisfied on the two bounding surfaces are different. However, the solution for this case can be reduced to that of case~(b) by considering odd (instead of even) solutions for $W$. For it is clear that an odd solution for $W$ satisfying the boundary conditions appropriate to case~(b) vanishes at $z = 0$; therefore, at $z = 0$ the boundary conditions on $W$ appropriate to a free surface are satisfied. Associated with an odd solution for $W$ is an even solution for $Z$; therefore, at $z = 0$, $DZ$ will vanish which is the further condition to be satisfied on a free surface. Consequently, a solution with $W$ odd and $Z$ even, suitable for case~(b) and applicable to a cell depth~$\frac{1}{2}d$, provides a solution for case~(c) applicable to a cell depth~$d$, and Rayleigh and Taylor numbers which are sixteen times smaller.

We consider then the odd solutions appropriate to case~(b) and obtain them once again by the variational method.

We now expand $F$ in a sine series of the form
\begin{equation}\label{eq:3-168b}
F = \sum_m A_m\sin 2m\pi z
\end{equation}
and follow the procedure outlined in \S\,(b) above. We find that the corresponding solutions for $W_m$ and $z_m$ are
\begin{equation}\label{eq:3-169}
W_m = c_{2m}\,\gamma_{2m}\sin 2m\pi z + \sum_{j=1}^{3} B_j^{(m)}\sinh q_j\, z
\end{equation}
and
\begin{equation}\label{eq:3-170}
Z_m = -\left(\frac{2\Omega}{\nu}\,d\right)\!\left(-2m\pi\gamma_{2m}\cos 2m\pi z + \sum_{j=1}^{3} B_j^{(m)}\frac{q_j}{x_j}\cosh q_j\, z\right),
\end{equation}
where $q_j$ and $x_j$ have the same meanings as in \S\,(b) and $c_{2m}$ and $\gamma_{2m}$ are defined as in equations~\eqref{eq:3-139} and~\eqref{eq:3-147} with $2m$ replacing $(2m+1)$. The constants of integration $B_j^{(m)}$ are similarly determined by the boundary conditions.

Substituting $F$ and $W$ in accordance with equations~\eqref{eq:3-168b} and~\eqref{eq:3-169} in the equation relating them (namely, equation~\eqref{eq:3-105}), we now have (cf.\ equation~\eqref{eq:3-156})
\begin{equation}\label{eq:3-171}
\sum_m A_m\, c_{2m}\sin 2m\pi z = Ra^2\sum_m A_m\!\left[c_{2m}\,\gamma_{2m}\sin 2m\pi z + \sum_{j=1}^{3} B_j^{(m)}\sinh q_j\, z\right].
\end{equation}

Multiplying this equation by $\sin 2n\pi z$ and integrating over the range of $z$, we obtain
\begin{equation}\label{eq:3-172}
\tfrac{1}{2}c_{2m}\, A_m = Ra^2\!\left[\tfrac{1}{2}c_{2m}\,\gamma_{2m}\, A_m + \sum_n (n|m)\,A_n\right],
\end{equation}
where
\begin{align}
(n|m) &= \sum_{j=1}^{3} B_j^{(m)} \int_{-1/2}^{+1/2} \sinh q_j\, z\sin 2m\pi z\,dz \notag\\
&= 4m\pi(-1)^n \sum_{j=1}^{3} \frac{B_j^{(m)}\sinh\tfrac{1}{2}q_j}{4n^2\pi^2 + q_j^2}, \label{eq:3-173}
\end{align}
and this leads to the characteristic equation
\begin{equation}\label{eq:3-174}
\left\|\tfrac{1}{2}c_{2m}\!\left[\frac{1}{Ra^2} - \gamma_{2m}\right]\!\delta_{mn} - (n|m)\right\| = 0.
\end{equation}

The explicit expression for the matrix element $(n|m)$ is found to be (cf.\ equation~\eqref{eq:3-165})
\begin{align}
(n|m) &= 8(-1)^{m+n+1}nm\pi^2\gamma_{2m}\,\gamma_{2n}\,\Phi^{(0)}\times \notag\\
&\quad\times\Bigl\{Y(c_{2m}+x)\bigl(1+\tfrac{c_{2m}}{x}\bigr)\coth\tfrac{1}{2}q + \notag\\
&\qquad + \tfrac{c_{2n}\,c_{2m}}{X^2+Y^2}\,\Phi^{(c)} - (c_{2n}+c_{2m})\,\Phi^{(c)} + \Phi^{(c)}\Bigr\}, \label{eq:3-175}
\end{align}
where $x$, $q$, $X$, $Y$ have the same meanings as in \S\,(b); and analogous to equations~\eqref{eq:3-166} and~\eqref{eq:3-167}, we now have
\begin{align}
\phi_1^{(c)} &= \frac{\alpha_2\sinh\alpha_1 + \alpha_2\sin\alpha_2}{\cosh\alpha_1 + \cos\alpha_2}, \qquad
\psi_1^{(c)} = \frac{\alpha_2\sinh\alpha_1 - \alpha_1\sin\alpha_2}{\cosh\alpha_1 + \cos\alpha_2}, \notag\\[6pt]
\phi_2^{(c)} &= \frac{\alpha_1\sinh\alpha_1 - \alpha_2\sin\alpha_2}{\cosh\alpha_1 - \cos\alpha_2}, \qquad
\psi_2^{(c)} = \frac{\alpha_2\sinh\alpha_1 + \alpha_1\sin\alpha_2}{\cosh\alpha_1 - \cos\alpha_2}, \label{eq:3-176}
\end{align}
\begin{align}
\Phi^{(c)} &= \frac{\tanh\tfrac{1}{2}q\,(\sinh^2\!\tfrac{1}{2}\alpha_1 + \sin^2\!\tfrac{1}{2}\alpha_2)(\cosh\alpha_1 + \cos\alpha_2)^{-1}}{\left\{q\frac{Y\phi^{(c)} - X\psi^{(c)}}{X^2+Y^2} + \frac{\psi^{(c)}}{z}\right\}} - \frac{\alpha_1+\alpha_2}{X^2+Y^2}\,Y\tanh\tfrac{1}{2}q. \notag
\end{align}

\begin{align}
\Phi_1^{(c)} &= (X^2 - Y^2 - xX)\psi_1^{(c)} - Y(2X-x)\phi_1^{(c)}, \notag\\
\Phi_2^{(c)} &= Y\phi_1^{(c)} - (X-x)\psi_1^{(c)}, \notag\\
\Phi_3^{(c)} &= (X^2 + Y^2 - xX)\psi_1^{(c)} - xY\phi_1^{(c)}. \label{eq:3-176b}
\end{align}

The secular equation~\eqref{eq:3-174} has been solved in various approximations. The results of such calculations are summarized in Table~IX and illustrated in Figs.~21 and~22 (see p.~96).

\begin{table}[ht]
\centering
\caption{Critical Rayleigh numbers and related constants for the case when one bounding surface is rigid and the other is free and the onset of instability is as stationary convection}
\label{tab:IX}
\begin{tabular}{cccccccc}
\hline
& & \multicolumn{3}{c}{$R_c$} & \multicolumn{1}{c}{Second} & \multicolumn{2}{c}{Third approximation} \\
\cline{3-5}\cline{7-8}
$\mathscr{T}$ & $a_c$ & First & Second & Third & approxi- & & \\
& & approx. & approx. & approx. & mation $A_1/A_0$ & $A_1/A_0$ & $A_2/A_1$ \\
\hline
$6{\cdot}25$ & $2{\cdot}68$ & $1{\cdot}120\times10^3$ & $1{\cdot}1085\times10^3$ & $+0{\cdot}0680$ & & \\
$3{\cdot}125\times10^1$ & $2{\cdot}70$ & $1{\cdot}148\times10^3$ & $1{\cdot}1395\times10^3$ & $+0{\cdot}06140$ & $+0{\cdot}06335$ & $-0{\cdot}00101$ \\
$1{\cdot}875\times10^2$ & $2{\cdot}975$ & $1{\cdot}393\times10^3$ & & $+0{\cdot}04935$ & & \\
$6{\cdot}250\times10^2$ & $3{\cdot}40$ & $2{\cdot}060\times10^3$ & $1{\cdot}6376\times10^4$ & $+0{\cdot}03570$ & $+0{\cdot}01559$ & $-0{\cdot}01198$ \\
$1{\cdot}875\times10^3$ & $4{\cdot}125$ & $3{\cdot}950\times10^3$ & & $+0{\cdot}01033$ & & \\
$6{\cdot}250\times10^3$ & $6{\cdot}005$ & $7{\cdot}291\times10^3$ & & $+0{\cdot}00604$ & & \\
$6{\cdot}250\times10^4$ & $7{\cdot}475$ & $1{\cdot}453\times10^4$ & $1{\cdot}4510\times10^5$ & $+0{\cdot}03702$ & $-0{\cdot}03701$ & $+0{\cdot}00491$ \\
$1{\cdot}875\times10^5$ & $9{\cdot}00$ & $2{\cdot}850\times10^4$ & & $+0{\cdot}01796$ & & \\
$6{\cdot}250\times10^5$ & $11{\cdot}05$ & $6{\cdot}177\times10^4$ & & $-0{\cdot}02774$ & & \\
& $13{\cdot}80$ & $1{\cdot}381\times10^5$ & $8{\cdot}2265\times10^5$ & $-0{\cdot}07634$ & $-0{\cdot}07907$ & $+0{\cdot}03261$ \\
$10^7$ & $26{\cdot}55$ & $1{\cdot}730\times10^5$ & $1{\cdot}7173\times10^6$ & $-0{\cdot}08461$ & $-0{\cdot}08504$ & $+0{\cdot}04698$ \\
$10^8$ & $58{\cdot}15$ & $3{\cdot}774\times10^5$ & $3{\cdot}7579\times10^6$ & $-0{\cdot}06739$ & $-0{\cdot}06841$ & $+0{\cdot}03876$ \\
\hline
\end{tabular}
\end{table}

\subsection*{(d) The origin of the $\mathscr{T}^{1/3}$- and the $\mathscr{T}^{2/3}$-laws}

From the results for the three sets of boundary conditions illustrated in Figs.~21 and~22, it is apparent that all three cases exhibit the same general features. This is particularly true of the asymptotic dependence on $\mathscr{T}$ of the critical Rayleigh number and the associated wave number. For the case of two free boundaries, the laws
\begin{equation}\label{eq:3-177}
R_c \to \text{constant}\ \mathscr{T}^{2/3} \quad \text{and} \quad a \to \text{constant}\ \mathscr{T}^{1/6}
\end{equation}
follow directly from the solution of the characteristic value problem. From the results of the calculations for the two other cases, it appears that the same power laws hold for them also though the constants of proportionality seem to depend slightly, but definitely, on the boundary conditions.

By going back to the original differential equations, we shall try to locate the common origin of the $\mathscr{T}^{1/3}$- and the $\mathscr{T}^{2/3}$-laws.

Consider the equation
\begin{equation}\label{eq:3-178}
(D^2 - a^2)^3 W + \mathscr{T}\,D^2 W = -Ra^2 W.
\end{equation}

When $\mathscr{T}\to\infty$, we expect that for $R$ in the neighbourhood of $R_c$, $a$ also will tend to infinity. But we do not expect that the solution for $W$, in this limit, will show any special behaviour which will make any of its higher derivatives become `disproportionately' large. The solution for the case of two free boundaries, as well as the progression of the values of the coefficients $A_m$ in the variational solutions for the other two cases (see Tables~VIII and~IX), support this latter expectation. Accordingly, keeping only the terms in $\mathscr{T}$ and in the highest power of $a$ in equation~\eqref{eq:3-178}, as relevant in the limit $\mathscr{T}\to\infty$, we have
\begin{equation}\label{eq:3-179}
\mathscr{T}\,D^2 W = -(Ra^2 - a^6)W.
\end{equation}

This is an equation of the second order in $W$. Consequently, we cannot satisfy all the boundary conditions of the problem; we can satisfy only two of the six boundary conditions. It would appear that the vanishing of $W$ on the boundaries is one condition which we must satisfy on all

accounts. The solution of equation~\eqref{eq:3-179} which vanishes for $z = 0$ and~1 is
\begin{equation}\label{eq:3-180}
W = A\sin n\pi z,
\end{equation}
where $n$ is an integer. The corresponding lowest characteristic value of $R$ is
\begin{equation}\label{eq:3-181}
R = \frac{1}{a^2}(a^6 + \pi^2\,\mathscr{T}).
\end{equation}

As a function of $a$, $R$ given by this equation attains its minimum when
\begin{equation}\label{eq:3-182}
4a^3 - \frac{2}{a^3}\pi^2\,\mathscr{T} = 0,
\end{equation}
or
\begin{equation}\label{eq:3-183}
a = (\tfrac{1}{2}\pi^2\,\mathscr{T})^{1/6};
\end{equation}
and the associated value of $R$ is
\begin{equation}\label{eq:3-184}
R = 3(\tfrac{1}{2}\pi^2\,\mathscr{T})^{2/3}.
\end{equation}

These results agree with those given in equation~\eqref{eq:3-133}.

In one sense, the foregoing arguments do not go much beyond the discussion for the case of two free boundaries. But it does suggest that the origin of the $\mathscr{T}^{1/3}$- and the $\mathscr{T}^{2/3}$-laws probably lies in the effective lowering of the order of equation~\eqref{eq:3-178} and a corresponding reduction in the number of boundary conditions that need be satisfied as $\mathscr{T}\to\infty$. It is worth pointing out that the preceding discussion is not `fine' enough to account for the differences in the constants of proportionality in the asymptotic relations which the variational calculations strongly suggest.

\section{The motions in the horizontal planes and the cell patterns at the onset of instability as stationary convection}\label{sec:3-28}

We now turn to a description of the motions in the horizontal plane and of the cell patterns which can emerge at the onset of instability as stationary convection.

In terms of the solutions for $w$ and $\zeta$, the components of the velocity in the horizontal plane are given by (equations~II~(110) and~(111))
\begin{equation}\label{eq:3-185}
\mathbf{u} = \frac{1}{a^2}\!\left(\frac{\partial^2 w}{\partial x\,\partial z}+\frac{\partial\zeta}{\partial y}\right), \quad
\mathbf{v} = \frac{1}{a^2}\!\left(\frac{\partial^2 w}{\partial y\,\partial z}-\frac{\partial\zeta}{\partial x}\right).
\end{equation}

To be specific, let
\begin{align}
w &= W(z)\cos a_x\, x\cos a_y\, y \label{eq:3-186}\\
\intertext{and}
\zeta &= Z(z)\cos a_x\, x\cos a_y\, y. \notag
\end{align}

Equations~\eqref{eq:3-185} then give
\begin{align}
\mathbf{u} &= -\frac{1}{a^2}(a_x\,DW\sin a_x\, x\cos a_y\, y + a_y\,Z\cos a_x\, x\sin a_y\, y), \notag\\
\mathbf{v} &= -\frac{1}{a^2}(a_y\,DW\cos a_x\, x\sin a_y\, y - a_x\,Z\sin a_x\, x\cos a_y\, y). \label{eq:3-187}
\end{align}

Since (cf.\ equation~\eqref{eq:3-122})
\begin{equation}\label{eq:3-188}
(D^2 - a^2)Z = -\left(\frac{2\Omega d}{\nu}\right)\!DW,
\end{equation}
we can also write
\begin{align}
\mathbf{u} &= \frac{1}{(2\Omega d)/a^2}\!\frac{1}{\nu}\!\bigl\{(a_y\sin a_x\, x\cos a_y\, y)(D^2 - a^2)Z - \notag\\
&\qquad\qquad - (a_x\cos a_x\, x\sin a_y\, y)\,Z\,\sqrt{\mathscr{T}}\bigr\}, \notag\\
\mathbf{v} &= \frac{1}{(2\Omega d)/a^2}\!\frac{1}{\nu}\!\bigl\{(a_y\cos a_x\, x\sin a_y\, y)(D^2 - a^2)Z + \notag\\
&\qquad\qquad + (a_x\sin a_x\, x\cos a_y\, y)\,Z\,\sqrt{\mathscr{T}}\bigr\}. \label{eq:3-189}
\end{align}

From these expressions for $\mathbf{u}$ and $\mathbf{v}$ we find
\begin{equation}\label{eq:3-190}
u^2 + v^2 = \frac{1}{(2\Omega d)^2}\!\frac{1}{\nu^2}(a_x^2\sin^2\!a_x\, x\cos^2\!a_y\, y + a_y^2\cos^2\!a_x\, x\sin^2\!a_y\, y)\times\bigl\{[(D^2 - a^2)Z]^2 + \mathscr{T}\,Z^2\bigr\}.
\end{equation}

Therefore, the average over a horizontal plane, we have
\begin{equation}\label{eq:3-191}
\langle u^2 + v^2\rangle = \frac{1}{4a^2}\frac{1}{(2\Omega d)^2}\bigl\{[(D^2 - a^2)Z]^2 + \mathscr{T}\,Z^2\bigr\}.
\end{equation}

For the case of two free boundaries, the solutions for $W$ and $Z$ appropriate for the lowest mode can be written as (cf.\ equation~\eqref{eq:3-126})
\begin{equation}\label{eq:3-192}
W = W_0\sin\pi z \quad \text{and} \quad Z = W_0\!\left(\frac{2\Omega d}{\nu}\right)\!\frac{\pi}{\pi^2 + a^2}\cos\pi z.
\end{equation}

A somewhat surprising result emerges when the foregoing solution for $Z$ is used in equation~\eqref{eq:3-191}. We find
\begin{equation}\label{eq:3-193}
\langle u^2 + v^2\rangle = \tfrac{1}{4}W_0^2\,\frac{\pi^2}{a^2(\pi^2 + a^2)}\bigl[(\pi^2 + a^2)^2 + \mathscr{T}\bigr]\cos^2\!\pi z.
\end{equation}

Letting $a^2 = \pi^2 x$ (as in \S\,27\,(a)), equation~\eqref{eq:3-129}), we have
\begin{equation}\label{eq:3-194}
\langle u^2 + v^2\rangle = \tfrac{1}{4}W_0^2\,\frac{(1+x)^2 + \mathscr{T}/\pi^4}{x(1+x)^2}\cos^2\!\pi z.
\end{equation}

On the other hand, the value of $x$ at the onset of instability is related to $\mathscr{T}$ by equation~\eqref{eq:3-131}; using this equation to eliminate $\mathscr{T}$ from~\eqref{eq:3-194}, we find that we are left with
\begin{equation}\label{eq:3-195}
\langle u^2 + v^2\rangle = \tfrac{1}{4}W_0^2\cos^2\!\pi z,
\end{equation}
a result which is independent of $\mathscr{T}$. Thus, \emph{for a layer of fluid confined between two free boundaries, the ratio, of the mean kinetic energies of the motions in the horizontal plane and in the vertical direction at the onset of stationary convection, is independent of rotation}.

For other boundary conditions, the invariance expressed by equation~\eqref{eq:3-195} is asymptotically true for $\mathscr{T}\to\infty$: for, according to the discussion in \S\,27\,(d), the general case tends to that for two free boundaries as $\mathscr{T}\to\infty$.

We now turn to the cell patterns. These have been described and beautifully illustrated by Veronis. The following account largely derives from his work.

\subsection*{(a) Rolls, rectangles, and squares}

The case of rectangular cells (of which rolls and squares are special cases) can be derived from the solutions for $\mathbf{u}$ and $\mathbf{v}$ given in~\eqref{eq:3-189}. For the case of two free boundaries (to which the present discussion on the cell patterns will be limited), we can insert for $Z$ its solution from~\eqref{eq:3-192}. We thus obtain the explicit formulae
\begin{align}
\mathbf{u} &= -\frac{\pi}{a^2}\!\left(a_x\sin a_x\, x\cos a_y\, y + \frac{\sqrt{\mathscr{T}}}{\pi^2 + a^2}\,a_y\cos a_x\, x\sin a_y\, y\right)\!W_0\cos\pi z, \notag\\[6pt]
\mathbf{v} &= -\frac{\pi}{a^2}\!\left(a_y\cos a_x\, x\sin a_y\, y - \frac{\sqrt{\mathscr{T}}}{\pi^2 + a^2}\,a_x\sin a_x\, x\cos a_y\, y\right)\!W_0\cos\pi z. \label{eq:3-196}
\end{align}

The corresponding solution for the vertical velocity is given by
\begin{equation}\label{eq:3-197}
w = W_0\cos a_x\, x\cos a_y\, y\sin\pi z.
\end{equation}

The case of rolls is of particular simplicity. The solution for this case can be obtained by setting $a_x = a$ and $a_y = 0$ in equations~\eqref{eq:3-196} and~\eqref{eq:3-197}; thus
\begin{align}
\mathbf{u} &= -\frac{\pi}{a(\pi^2 + a^2)}W_0\sin ax\cos\pi z, \label{eq:3-198}\\
\mathbf{v} &= +\frac{\sqrt{\mathscr{T}}}{a(\pi^2 + a^2)}W_0\sin ax\cos\pi z, \label{eq:3-199}
\end{align}
and
\begin{equation}
w = W_0\cos ax\sin\pi z. \notag
\end{equation}

When $\mathscr{T} = 0$, there are no motions in the $y$-direction and the streamlines are confined to planes normal to the direction of the rolls. When the system is rotated, the Coriolis acceleration induces longitudinal motions. Since in the present case
\begin{equation}\label{eq:3-200}
\frac{v}{u} = -\frac{\sqrt{\mathscr{T}}}{\pi^2 + a^2} = \text{constant},
\end{equation}
the streamlines are again confined to planes; but these planes are inclined to the $z$-axis (see Fig.~23\,$b,c$). Under these circumstances we can define a wave number $a_s$ which will describe the motions in the oblique planes containing the streamlines; it is given by
\begin{equation}\label{eq:3-201}
a_s = a\cos(\tan^{-1} v/u) = \frac{a}{\sqrt{1+v^2/u^2}}.
\end{equation}

Substituting for $v/u$ from~\eqref{eq:3-200}, we obtain
\begin{equation}\label{eq:3-202}
a_s^2 = \frac{a^2(\pi^2 + a^2)^2}{(\pi^2 + a^2)^2 + \mathscr{T}}.
\end{equation}

\figurePlaceholder{23}{(a) A top view of two-dimensional rolls in a non-rotating fluid. (b) The same view in a system rotating counterclockwise. (c) A perspective view of the particle motions in a roll in the rotating case. The arrows indicate the direction of particle motions.}

The quantity on the right-hand side of this equation is, apart from a numerical factor, the reciprocal of the expression which we have already shown to be independent of $\mathscr{T}$ (cf.\ the arguments leading from equation~\eqref{eq:3-193} to~\eqref{eq:3-195}); it has the value $\frac{1}{4}\pi^2$. Thus,
\begin{equation}\label{eq:3-203}
a_s^2 = \tfrac{1}{4}\pi^2;
\end{equation}
but this is the same as the formula which gives the wave number of the roll in the absence of rotation (see equation~II~(195)). Hence, \emph{the wavelength of the roll measured in the plane containing the streamlines is independent of rotation}. This remarkable result is due to Veronis.

Considering next the case of square cells, we obtain the corresponding solutions by setting $a_x = a_y = a/\sqrt{2}$ in~\eqref{eq:3-196}. We obtain
\begin{equation}\label{eq:3-204}
\mathbf{u} = -\frac{\pi}{a\sqrt{2}}\!\left(\sin\frac{ax}{\sqrt{2}}\cos\frac{ay}{\sqrt{2}} + \frac{\sqrt{\mathscr{T}}}{\pi^2 + a^2}\cos\frac{ax}{\sqrt{2}}\sin\frac{ay}{\sqrt{2}}\right)\!W_0\cos\pi z,
\end{equation}
and similarly for $\mathbf{v}$.

The streamlines in the absence of rotation in such square cells have been

\figurePlaceholder{24a}{A square cell in a rotating fluid. The particle motions are from the centre outward along spiral curves. The dashed curves form the boundary of the square cell.}

\noindent illustrated in Chapter~II (Fig.~6, p.~46). The dashed curves in Fig.~24\,$a$ correspond to the sides of the square in Fig.~6. In the presence of rotation, the Coriolis acceleration causes the fluid to turn as it moves towards or outwards from the centre; motions along the curves of constant $w$, as well as transverse to them, occur. Fig.~24\,$b$ is a perspective of the complete trajectory as sketched by Veronis. The drawing illustrates how a fluid element starting at the centre of a cell spirals upwards in a clockwise direction (for a counter-clockwise rotation $\Omega$), and as it approaches the top it crosses over towards a corner of a cell and starts

\figurePlaceholder{24b}{A perspective sketch of the path of a fluid particle in a square cell.}

\noindent spiralling downwards in a counter-clockwise direction. When it arrives at the mid-plane ($z = \frac{1}{2}$), it reverses its sense of spiralling and proceeds downwards, and as it reaches the bottom it crosses back towards the centre of the cell, and starts spiralling upwards towards the centre in a clockwise direction. The reversal of the rotation at the mid-plane is associated with the circumstance that here $DW$ and, therefore, also, $\nabla_1\cdot\mathbf{u}_1$ change sign.

The square cells shown in Fig.~24\,$b$ correspond to the basic geometry of the vertical velocity. The distortion of the cell consequent on the rotation is not shown.

\subsection*{(b) Hexagons}

Christopherson's solution for a hexagonal pattern and appropriate for two free boundaries is (cf.\ equation~II~(248))
\begin{equation}\label{eq:3-205}
w = \frac{1}{3}\!\left(2\cos\frac{2\pi}{L\sqrt{3}}\,x\cos\frac{2\pi}{3L}\,y + \cos\frac{4\pi}{3L}\,y\right)\!W_0\sin\pi z.
\end{equation}

The corresponding solution for the vertical component of the vorticity is
\begin{equation}\label{eq:3-206}
\zeta = \frac{1}{3}\!\left(\frac{2\Omega d}{\nu}\right)\!\frac{\pi}{\pi^2 + a^2}\!\left(2\cos\frac{2\pi}{L\sqrt{3}}\,x\cos\frac{2\pi}{3L}\,y + \cos\frac{4\pi}{3L}\,y\right)\!W_0\cos\pi z.
\end{equation}

Substituting for $w$ and $\zeta$ in accordance with these solutions in~\eqref{eq:3-187} we obtain
\begin{align}
\mathbf{u} &= -\frac{\pi}{3a^2}\!\left(\frac{4\pi}{L\sqrt{3}}\sin\frac{2\pi}{L\sqrt{3}}\,x\cos\frac{2\pi}{3L}\,y + \right.\notag\\
&\qquad\left. +\frac{4\pi}{3L}\frac{\sqrt{\mathscr{T}}}{\pi^2 + a^2}\!\left(\cos\frac{2\pi}{L\sqrt{3}}\,x + 2\cos\frac{2\pi}{3L}\,y\right)\!\sin\frac{2\pi}{3L}\,y\right)\!W_0\cos\pi z, \notag\\[8pt]
\mathbf{v} &= -\frac{\pi}{3a^2}\!\left(\frac{4\pi}{3L}\!\left(\cos\frac{2\pi}{L\sqrt{3}}\,x + 2\cos\frac{2\pi}{3L}\,y\right)\!\sin\frac{2\pi}{3L}\,y - \right.\notag\\
&\qquad\left. -\frac{4\pi}{L\sqrt{3}}\frac{\sqrt{\mathscr{T}}}{\pi^2 + a^2}\sin\frac{2\pi}{L\sqrt{3}}\,x\cos\frac{2\pi}{3L}\,y\right)\!W_0\cos\pi z. \label{eq:3-207}
\end{align}

\figurePlaceholder{25a}{A top view of seven rotating hexagonal cells. The fluid particles follow spiral paths from the centre toward the corners. The dashed lines form the boundaries of the centre cell.}

Fig.~25\,$a$ shows the streamlines as they would appear at the top. This pattern should be contrasted with that illustrated in Fig.~7\,$b$ (p.~50) for a non-rotating system. The dashed spirals form the boundary of the central cell. It will be observed that the source at the centre of the cell provides for the flow downwards at the six corners; and each corner

\figurePlaceholder{25b}{A perspective sketch of the path of a fluid particle in a hexagonal cell.}

\noindent receives fluid from three adjacent cells. In this last respect it is not different from the non-rotating case.

A perspective drawing of the complete trajectory is shown in Fig.~25\,$b$. And finally, in Fig.~26 we illustrate the manner in which the cell is distorted by the rotation.

\figurePlaceholder{26}{A perspective sketch of a hexagonal cell as it is distorted by the rotation of the fluid.}

\subsection*{(c) The limiting behaviour of the streamlines for $\mathscr{T}\to\infty$}

It is clear from the foregoing illustrations of the cell patterns that as $\mathscr{T}$ increases, the streamlines become increasingly close wound spirals. In the limit $\mathscr{T}\to\infty$, when convection in fact ceases, the streamlines become closed curves. Thus, for $\mathscr{T}$ sufficiently large, the dominant terms in the solution~\eqref{eq:3-196} for $\mathbf{u}$ and $\mathbf{v}$ are, for example,
\begin{align}
\mathbf{u} &\to -\frac{\pi\sqrt{\mathscr{T}}}{a^2(\pi^2+a^2)}\,a_y\,W_0\cos a_x\, x\sin a_y\, y\cos\pi z, \notag\\[4pt]
\mathbf{v} &\to +\frac{\pi\sqrt{\mathscr{T}}}{a^2(\pi^2+a^2)}\,a_x\,W_0\sin a_x\, x\cos a_y\, y\cos\pi z. \label{eq:3-208}
\end{align}

The corresponding equation for the streamlines is
\begin{equation}\label{eq:3-209}
\cos a_x\, x\cos a_y\, y = \text{constant}
\end{equation}
or,
\begin{equation}\label{eq:3-210}
w = \text{constant} \quad (\text{for } z = \text{constant}).
\end{equation}

This last result is not restricted to rectangular cells: it applies equally to all cell patterns.

Whereas, in the absence of rotation $\mathbf{u}_\perp$ is parallel to $\nabla_1 w$, in the limit $\mathscr{T}\to\infty$ the streamlines tend to coincide with the curves $w = \text{constant}$. For $\mathscr{T}$ large but finite, we may describe the motions in the horizontal plane as consisting, principally, of motions around the curves of constant $w$; and superposed on this is a slow outward, or inward, directed radial motion which induces a spiral pattern. Near the centres of the cells, the spiral patterns tend to become equiangular. It is the small radial motion in the direction of $\nabla_1 w$ that is responsible for all phenomena associated with convection when $\mathscr{T}$ becomes large.

% ========================================================================
% [MERGED] Content below was originally in chapter_3_part2.tex
% ========================================================================

%% --- Section 29 ---
\section{On the onset of convection as overstability. The solution for the case of two free boundaries}\label{sec:3-29}

We now take up the question, set aside in \S~26, whether or not instability can arise as oscillations of increasing amplitude, i.e., as overstability. This requires us to return to the general equations~\eqref{eq:3-93}--\eqref{eq:3-95} which include the time constant~$\sigma$. In \S~30 we shall describe how best one may treat these equations to decide under what conditions stationary or overstable convection arise. In this section we shall restrict ourselves to the case of two free boundaries; for in this case the problem can be solved by elementary methods and it suggests how one may treat the general case.

By applying the operator $(D^2 - a^2 - \sigma)(D^2 - a^2 - p\sigma)$ to equation~\eqref{eq:3-95}, we can eliminate $Z$ and $\Theta$ and obtain
\begin{equation}\label{eq:3-211}
(D^2 - a^2 - p\sigma)(D^2 - a^2 - \sigma)^2(D^2 - a^2) + \mathscr{T} D^2] W = -Ra^2(D^2 - a^2 - \sigma)W.
\end{equation}

As in \S~27\,$(a)$, we can show that in this case also the proper solution for $W$ belonging to the lowest mode is
\begin{equation}\label{eq:3-212}
W = W_0 \sin \pi z.
\end{equation}

Substituting this solution for $W$ in equation~\eqref{eq:3-211}, we obtain the characteristic equation
\begin{equation}\label{eq:3-213}
(\pi^2 + a^2 + p\sigma)[(\pi^2 + a^2 + \sigma)^2(\pi^2 + a^2) + \pi^2\mathscr{T}] = Ra^2(\pi^2 + a^2 + \sigma),
\end{equation}
where it must be remembered that $\sigma$ can be complex. Letting
\begin{equation}\label{eq:3-214}
x = \frac{a^2}{\pi^2}, \quad i\sigma_1 = \frac{\sigma}{\pi^2}, \quad R_1 = \frac{R}{\pi^4}, \quad \text{and} \quad T_1 = \frac{\mathscr{T}}{\pi^4},
\end{equation}
we can rewrite equation~\eqref{eq:3-213} in the form
\begin{equation}\label{eq:3-215}
R_1 = \frac{1}{x}(1 + x + ip\sigma_1)(1 + x)(1 + x + i\sigma_1) + \frac{T_1}{1 + x + i\sigma_1}\bigg\}.
\end{equation}

It is apparent from equation~\eqref{eq:3-215} that for an arbitrarily assigned $\sigma_1$, $R_1$ will be complex. But the physical meaning of $R_1$ requires it to be real. Consequently, the condition that $R_1$ be real implies a relation between the real and imaginary parts of $\sigma_1$. Since our principal interest is to specify the critical Rayleigh number for the onset of instability via a state of purely oscillatory motions, we shall in the first instance suppose that $\sigma_1$ in equation~\eqref{eq:3-215} is real and seek the conditions for such solutions to exist. This will suffice to answer the principal question as to when instability will set in as stationary convection and when as overstable oscillations. To establish that the criteria so obtained are both necessary and sufficient, we must consider the characteristic equation~\eqref{eq:3-215} more generally, and this we will do in \S~$(a)$ below.

We shall assume, then, that in equation~\eqref{eq:3-215} $\sigma_1$ is real. On collecting the real and imaginary parts of the expression on the right-hand side of equation~\eqref{eq:3-215}, we have
\begin{equation}\label{eq:3-216}
R_1 = \frac{1+x}{x}\bigg\{(1+x)^2 - p\sigma_1^2 + \frac{T_1}{1+x}\,\frac{(1+x)^2 + p\sigma_1^2}{(1+x)^2 + \sigma_1^2} + i\sigma_1\bigg[(1+x)(1+p) - T_1\,\frac{(1-p)}{(1+x)^2 + \sigma_1^2}\bigg]\bigg\}.
\end{equation}

The real and the imaginary parts of this equation must vanish separately. We thus obtain the pair of equations:
\begin{equation}\label{eq:3-217}
R_1 = \frac{1+x}{x}\bigg\{(1+x)^2 - p\sigma_1^2 + \frac{T_1}{1+x}\,\frac{(1+x)^2 + p\sigma_1^2}{(1+x)^2 + \sigma_1^2}\bigg\}
\end{equation}
and
\begin{equation}\label{eq:3-218}
(1+x)(1+p) = T_1\,\frac{(1-p)}{(1+x)^2 + \sigma_1^2}.
\end{equation}

A relation which follows from equation~\eqref{eq:3-218} is
\begin{equation}\label{eq:3-219}
\frac{T_1}{1+x}\,\frac{(1+x)^2 + p\sigma_1^2}{(1+x)^2 + \sigma_1^2} = \frac{T_1}{1+x}\,\frac{T_1(1-p)\sigma_1^2}{(1+x)^2 + \sigma_1^2} = \frac{T_1}{1+x} - (1+p)\sigma_1^2.
\end{equation}

Using this relation in~\eqref{eq:3-217}, we have
\begin{equation}\label{eq:3-220}
R_1 = \frac{1}{x}\big[(1+x)^3 + T_1 - (1+x)(1+2p)\sigma_1^2\big].
\end{equation}

Next, solving equation~\eqref{eq:3-218} for $\sigma_1^2$, we find
\begin{equation}\label{eq:3-221}
\sigma_1^2 = \frac{T_1}{1+x}\,\frac{1-p}{1+p} - (1+x)^2.
\end{equation}

Using this expression for $\sigma_1^2$ in~\eqref{eq:3-220}, we obtain on further simplification the result
\begin{equation}\label{eq:3-222}
R_1 = 2(1+p)\frac{1}{x}\bigg[(1+x)^3 + \frac{p^2}{(1+p)^2}\,T_1\bigg].
\end{equation}

Equations~\eqref{eq:3-221} and~\eqref{eq:3-222} are the equations which must be satisfied if overstability is to occur for a wave number corresponding to $x$ and a Taylor number corresponding to $T_1$.

One conclusion which equation~\eqref{eq:3-221} enables us to draw, at once, is that solutions describing overstability cannot occur if
\begin{equation}\label{eq:3-223}
\frac{T_1}{1+x}\,\frac{1-p}{1+p} < 1;
\end{equation}
for $\sigma_1^2$ would then be negative, contrary to hypothesis. Clearly, \eqref{eq:3-223} will, \emph{a fortiori}, be the case if $p > 1$. Accordingly, \emph{for $p > 1$, overstability cannot occur and the principle of the exchange of stabilities is valid}.

For overstability to be at all possible, $p$ must be less than one; and even when this is the case, we shall obtain real frequencies, $\sigma_1$, only if
\begin{equation}\label{eq:3-224}
T_1 > \frac{1+p}{1-p}(1+x)^3.
\end{equation}

For a given $T_1$, overstable solutions are therefore possible only for $x < x_a$ where $x_a$ is such that
\begin{equation}\label{eq:3-225}
(1+x_a)^3 = T_1\,\frac{1-p}{1+p}.
\end{equation}

When $x = x_a$ and $\sigma_1^2 = 0$ and (cf.\ equation~\eqref{eq:3-220})
\begin{equation}\label{eq:3-226}
R_1 = \frac{1}{x_a}\big[(1+x_a)^3 + T_1\big];
\end{equation}
and this is, as one should expect, the value of $R_1$ at which stationary convection will occur for a wave number corresponding to $x_a$ (see equation~\eqref{eq:3-130}). For $x > x_a$, overstability is not possible for the given $p$ and $T_1$; and the onset of instability as stationary convection remains the only possibility. For $x < x_a$, overstability is possible; and the question of discriminating between the two manners of instability arises. In making the discrimination, we shall suppose, other things being equal, that manner of instability will occur that allows a solution for a lower Rayleigh number. While this is intuitively obvious, it requires proof; this is furnished in \S~$(a)$ below.

Consider the $(x, R_1)$-plane (see Fig.~27). In this plane, we have first the curve,
\begin{equation}\label{eq:3-227}
R_1^{(c)} = \frac{1}{x}\big[(1+x)^3 + T_1\big],
\end{equation}
which defines the locus of states which are marginal with respect to stationary convection. The overstable solutions branch off from this

\figurePlaceholder{27}{The disposition of the curves for marginal stability in the $(x, R_1)$-plane for a Taylor number, $T = 10^6$, and for various values of the Prandtl number~$p$. The curve labelled `convection' applies for all values of~$p$ and defines the locus $R_1^{(c)}$ (equation~227). The remaining curves are the overstable loci $R_1^{(o)}$ for values of~$p$ by which they are labelled. The minimum for the curve $C$ ($p = 0{\cdot}6120$) and for the convection curve occurs for the same value of~$R_1$.}

\noindent locus at the point $x_a$ (see equation~\eqref{eq:3-226}); and for $x < x_a$ they are described by (cf.\ equations~\eqref{eq:3-220} and~\eqref{eq:3-222})
\begin{equation}\label{eq:3-228}
R_1^{(o)} = R_1^{(c)} - (1+2p)\sigma_1^2\,\frac{1+x}{x} = 2(1+p)\frac{1}{x}\bigg[(1+x)^3 + \frac{p^2}{(1+p)^2}\,T_1\bigg],
\end{equation}
where the first form shows that $R_1^{(o)}$ (when this branch of the solution exists) is \emph{always} less than $R_1^{(c)}$. Depending on $p$ and $T_1$, the branch point $x_a$ can occur either before, or after, the point, $x_{\min}^{(c)}$, at which $R_1^{(c)}$ attains its minimum. If $x_a > x_{\min}^{(c)}$, then it is clear that for all $x < x_a$, overstability is the preferred manner of instability: these are shown by the different curves labelled by $B$, $C$, $D$, and $E$ in Fig.~27. It is apparent from this figure that the case which will distinguish whether, with increasing $R_1$ (for a given $T_1$), overstability or stationary convection will first manifest itself, is the one for which the minima of the two curves $R_1^{(c)}$ and $R_1^{(o)}$ are equal. Thus, if the branch point $x_a$ occurs for a value of $x$ less than what corresponds to $C$ in Fig.~27, we can exclude overstability; on the other hand, if the branch point occurs for a value of $x$ greater than what corresponds to $C$, we can exclude stationary convection.

We shall now show that there exists a value of $p$ such that $R_{1,\min}^{(c)}$ approaches $R_{1,\min}^{(o)}$ from above as $T_1 \to \infty$. For according to equations~\eqref{eq:3-227} and~\eqref{eq:3-228}, the asymptotic behaviours of $R_{1,\min}^{(c)}$ and $R_{1,\min}^{(o)}$ are given by (cf.\ equation~\eqref{eq:3-139})
\begin{equation}\label{eq:3-229}
R_{1,\min}^{(c)} \to 3(\tfrac{4}{3}T_1)^{\frac{2}{3}},
\end{equation}
\begin{equation}\label{eq:3-230}
R_{1,\min}^{(o)} \to 2(1+p)\bigg\{\tfrac{3}{2}\bigg[\frac{p^2}{2\,(1+p)^2}\,T_1\bigg]^{\frac{2}{3}}\bigg\}.
\end{equation}

The condition that $R_{1,\min}^{(c)} \to R_{1,\min}^{(o)}$ as $T_1 \to \infty$, clearly, requires that
\begin{equation}\label{eq:3-231}
2(1+p)\bigg[\frac{p^2}{(1+p)^2}\bigg]^{\frac{2}{3}} = 1,
\end{equation}
or
\begin{equation}\label{eq:3-232}
\frac{p^{\frac{4}{3}}}{(1+p)^{\frac{1}{3}}} = 1.
\end{equation}

It is found that the required root of this equation is
\begin{equation}\label{eq:3-233}
p = 0{\cdot}67659 = p^* \;\text{(say)}.
\end{equation}

From the monotone dependence of the minimum of the curve,
\begin{equation}\label{eq:3-234a}
y = \frac{1}{x}\big[(1+x)^3 + X\big],
\end{equation}
on $X$, we conclude that for $p > p^*$, $R_{1,\min}^{(c)} < R_{1,\min}^{(o)}$ for all $T_1$. Hence, for $1 > p > p^*$, the situations indicated by $D$ or $E$ will prevail for all $T_1$. Therefore, \emph{for $p > p^*$, instability will always manifest itself, first, as stationary convection}.

For $p < p^*$, the situation indicated by $C$ occurs for a finite $T_1$, say $T_1^{(p)}$. For $T_1 < T_1^{(p)}$, the disposition of the curves $R_1^{(o)}$, with respect to $R_1^{(c)}$, is as shown by the curves $D$ and $E$ in Fig.~27; when $T_1 = T_1^{(p)}$, the situation is as shown by the curve $C$; and for $T_1 > T_1^{(p)}$ the disposition of the curves $R_1^{(o)}$, with respect to $R_1^{(c)}$, is as shown by the curves $B$ and $A$. We conclude, then, that \emph{for $p < p^*$ there exists a $T_1^{(p)}$ such that for $T_1 < T_1^{(p)}$, the onset of instability will be as stationary convection, while for $T_1 > T_1^{(p)}$ it will be as overstability}.

There is no simple formula which gives $T_1^{(p)}$ as a function of $p$: it is simply determined by the condition that $R_{1,\min}^{(c)}$ and $R_{1,\min}^{(o)}$ are equal for

\begin{table}[ht]
\centering
\caption{Taylor numbers above which instability occurs as overstability}
\label{tab:X}
\begin{tabular}{cccccc}
\hline
$p$ & $T^{(p)}$ & $R$ & $p$ & $T^{(p)}$ & $R$ \\
\hline
$0$ & & & $0{\cdot}55$ & $15{,}370$ & 7748 \\
$0{\cdot}1$ & 548 & 1315 & $0{\cdot}60$ & $68{,}150$ & $16{,}290$ \\
$0{\cdot}2$ & 728 & 1431 & $0{\cdot}63$ & $2{\cdot}588\times 10^5$ & $3{\cdot}870\times 10^4$ \\
$0{\cdot}3$ & 990 & 1669 & $0{\cdot}65$ & $1{\cdot}233\times 10^6$ & $1{\cdot}059\times 10^5$ \\
$0{\cdot}4$ & 3103 & 2890 & $0{\cdot}6766$ & $\infty$ & $\infty$ \\
$0{\cdot}5$ & 8595 & 4970 & & & \\
\hline
\end{tabular}
\end{table}

\figurePlaceholder{28}{The variation of $T^{(p)}$ with the Prandtl number~$p$.}

\noindent $T_1 = T_1^{(p)}$. Values of $T^{(p)}$ determined by this condition are given in Table~X for a few values of $p$; and Fig.~28 shows the division of the $(p, T)$-plane into regions where the manner of instability at onset is one or the other; the part of this plane where overstability can occur for Rayleigh numbers higher than are required for the onset of instability is bounded by $p = 1$ and the locus~$T^{(p)}$.

For $p < p^*$, the complete $(R_c, T)$-relation can be deduced from the results given in Table~VII by a simple transformation of scales: for, according to equation~\eqref{eq:3-222}, we have only to interpret $T$ in Table~VII

\begin{table}[ht]
\centering
\caption{Critical Rayleigh numbers and related constants for the onset of overstability in case both bounding surfaces are free and $p = 0{\cdot}025$}
\label{tab:XI}
\small
\begin{tabular}{ccccccc}
\hline
$T$ & $a_c$ & $R_c$ & $p|\sigma|$ & & & \\
\hline
0 & $2{\cdot}33$ & & $1{\cdot}348\times 10^4$ & & \text{imaginary} & \\
$1{\cdot}681\times 10^3$ & $2{\cdot}70$ & $1{\cdot}014\times 10^4$ & $1{\cdot}388\times 10^5$ & & $1{\cdot}504$ & \\
$1{\cdot}681\times 10^4$ & $2{\cdot}944$ & $3{\cdot}079\times 10^4$ & $1{\cdot}064\times 10^6$ & & $1{\cdot}262$ & \\
$8{\cdot}493\times 10^4$ & $3{\cdot}278$ & $6{\cdot}184\times 10^4$ & $2{\cdot}615\times 10^6$ & & $1{\cdot}349$ & \\
$1{\cdot}875\times 10^5$ & $3{\cdot}710$ & $8{\cdot}055\times 10^4$ & $3{\cdot}435\times 10^6$ & & $1{\cdot}266$ & \\
$3{\cdot}364\times 10^5$ & $4{\cdot}220$ & $1{\cdot}007\times 10^5$ & $4{\cdot}713\times 10^6$ & & $1{\cdot}164$ & \\
$8{\cdot}495\times 10^5$ & $5{\cdot}011$ & $1{\cdot}303\times 10^5$ & $7{\cdot}513\times 10^6$ & & $1{\cdot}037$ & \\
$1{\cdot}681\times 10^6$ & $5{\cdot}968$ & $1{\cdot}937\times 10^5$ & $1{\cdot}183\times 10^7$ & & $0{\cdot}944$ & \\
$3{\cdot}943\times 10^6$ & $6{\cdot}961$ & $2{\cdot}851\times 10^5$ & $2{\cdot}092\times 10^7$ & & $0{\cdot}5029$ & \\
$1{\cdot}681\times 10^7$ & $8{\cdot}606$ & $4{\cdot}330\times 10^5$ & $4{\cdot}858\times 10^7$ & & $0{\cdot}5608$ & \\
$3{\cdot}943\times 10^7$ & $10{\cdot}43$ & $6{\cdot}360\times 10^5$ & $7{\cdot}728\times 10^7$ & & $0{\cdot}5618$ & \\
$1{\cdot}681\times 10^8$ & $12{\cdot}86$ & $9{\cdot}497\times 10^5$ & $1{\cdot}501\times 10^8$ & & $0{\cdot}4931$ & \\
$3{\cdot}943\times 10^8$ & $13{\cdot}06$ & $1{\cdot}370\times 10^6$ & $1{\cdot}784\times 10^8$ & & $0{\cdot}4790$ & \\
$1{\cdot}681\times 10^9$ & $17{\cdot}05$ & $2{\cdot}067\times 10^6$ & $3{\cdot}447\times 10^8$ & & $0{\cdot}3896$ & \\
$1{\cdot}681\times 10^{10}$ & $28{\cdot}02$ & $4{\cdot}470\times 10^6$ & $8{\cdot}958\times 10^8$ & & $0{\cdot}2175$ & \\
$1{\cdot}681\times 10^{11}$ & $38{\cdot}87$ & $8{\cdot}470\times 10^6$ & $8{\cdot}476\times 10^9$ & & $0{\cdot}0320$ & \\
$1{\cdot}681\times 10^{12}$ & $191{\cdot}51$ & $2{\cdot}075\times 10^7$ & & & & \\
\hline
\end{tabular}
\end{table}

\noindent as now signifying $p^4T/(1+p)^2$ and multiply the values under $R_c$ by $2(1+p)$. The $(R_c, T)$-relations derived in this manner are illustrated in Fig.~29. For later comparisons with results derived for other boundary conditions, the numerical form of the relation for $p = 0{\cdot}025$ is given in Table~XI.

It will be observed that in accordance with the result given in~\eqref{eq:3-229}, along the various overstable solutions $R \to \text{constant}\cdot T^{\frac{2}{3}}$ as $T \to \infty$. The explicit forms of the asymptotic behaviours of the Rayleigh number and the wave number of the associated disturbance for the onset of overstability may be noted; they are:
\begin{equation}\label{eq:3-234}
R_c^{(o)} \to \frac{6p^{\frac{4}{3}}}{(1+p)^{\frac{1}{3}}}\,(\tfrac{1}{2}\pi^2 T)^{\frac{2}{3}} \quad (p^2 T \to \infty).
\end{equation}

\begin{equation}\label{eq:3-235}
a_c^{(o)} \to \bigg(\frac{p}{1+p}\bigg)^{\frac{1}{3}} (\tfrac{1}{2}\pi^2 T)^{\frac{1}{6}} \quad (p^2 T \to \infty).
\end{equation}

The corresponding behaviour of $|\sigma| = \pi^2\sigma_1$ can be deduced from equation~\eqref{eq:3-221}. We find
\begin{equation}\label{eq:3-236}
|\sigma| \to \frac{(2-3p)^{\frac{1}{2}}}{|p(1+p)|^{\frac{1}{2}}}\,(\tfrac{1}{2}\pi^2 T)^{\frac{1}{2}} \quad (p^2 T \to \infty),
\end{equation}
where it may be noted that according to the definition of $\sigma$ (see equation~\eqref{eq:3-92}),
\begin{equation}\label{eq:3-237}
|p| = \frac{2|\sigma|}{\Omega} = \frac{2|\sigma_1|}{\sqrt{T_1}}.
\end{equation}

A further relation which follows from the foregoing is
\begin{equation}\label{eq:3-238}
|p|a_c^{(o)} \to 2\pi\Omega^{\frac{1}{2}}\,\frac{(1-1{\cdot}5p^{\frac{1}{2}})}{1+p} \quad (p^2 T \to \infty).
\end{equation}

It should be noted that the asymptotic relations~\eqref{eq:3-234}, \eqref{eq:3-236}, and~\eqref{eq:3-238} are valid only for $p^2 T \to \infty$, \emph{not} for $T \to \infty$, a distinction which is important when $p \to 0$ (see \S~32).

\figurePlaceholder{29}{The $(R_c, T)$-relations for a rotating horizontal layer of fluid heated below. The curves have been derived for the case when both bounding surfaces are free. The curve labelled `convection curve' is the $(R_c, T)$-relation for the onset of ordinary cellular convection. The remaining curves are the corresponding relations for the onset of overstability. The values of~$p$ to which the various curves refer are shown at the top of each curve. It will be seen that for each value of $p < 0{\cdot}677$, the instability sets in as ordinary cellular convection for $T$ less than a certain $T^{(p)}$ while it sets in as overstability for $T > T^{(p)}$.}

\subsection*{(a) The nature of the roots of the characteristic equation~\eqref{eq:3-215}}

The preceding discussion of the dependence of the manner of instability on $p$ and $T$ was carried out in terms of the explicit solutions for two special cases of equation~\eqref{eq:3-215}: the case when $\sigma_1 = 0$ and the marginal state is a stationary one; and the case when $\sigma_1$ is real and the marginal state is an oscillatory one. It was further assumed that the manner of instability which will occur at onset is the one which allows a solution for a lower Rayleigh number. It was pointed out that this assumption requires justification; and this we shall now provide.

In considering the roots of equation~\eqref{eq:3-215}, generally, we shall find it convenient to replace $i\sigma_1$ by $\sigma_1$ so that overstability now corresponds

\noindent to $\sigma_1$ being a pure imaginary number. With this replacement, the equation we have to consider may be written as
\begin{equation}\label{eq:3-239}
(1+x+p\sigma_1)\big[(1+x+\sigma_1)^2(1+x) + \frac{T_1}{1+x}\,d^2\big] = R_1(1+x+\sigma_1)\,\frac{x}{1+x},
\end{equation}
where
\begin{equation}\label{eq:3-240}
\begin{aligned}
B &= \frac{1}{p}(1+x)(1+2p), \\[4pt]
C &= \frac{1}{p}\big[(2+p)(1+x)^2 + p\,\frac{T_1}{1+x} - R_1\,\frac{x}{1+x}\big], \\[4pt]
D &= \frac{1}{p}\big[(1+x)^3 + T_1 - R_1\,x\big].
\end{aligned}
\end{equation}

This is a cubic equation for $\sigma_1$; its explicit form is
\begin{equation}\label{eq:3-241}
\sigma_1^3 + B\sigma_1^2 + C\sigma_1 + D = 0,
\end{equation}

A combination of the coefficients of the cubic~\eqref{eq:3-241}), which is important for the characterization of its roots, is
\begin{equation}\label{eq:3-242}
BC - D = \frac{1+p}{p}\bigg[(2+p)\bigg[(1+x)^2 + T_1\,\frac{p^2}{(1+p)^2}\bigg] - R_1\,x\bigg].
\end{equation}

We may first observe that $D = 0$ and $BC - D = 0$ define the solutions which we have denoted by $R_1^{(c)}$ and $R_1^{(o)}$, respectively (see equations~\eqref{eq:3-227} and~\eqref{eq:3-228}). It will be recalled that these solutions give the Rayleigh numbers for the onset of stationary convection and overstability. This is consistent with equation~\eqref{eq:3-241}: for $\sigma_1 = 0$ is clearly a root of the equation if $D = 0$; while, if $\sigma_1$ should be purely imaginary, then on equating the real and imaginary parts of equation~\eqref{eq:3-241} separately, we should obtain $|\sigma_1|^2 = C = D/B$.

Now when $R_1 = 0$, $C$, $D$, and $BC - D$, are all positive; and $B$ is, of course, always positive. When $D > 0$, equation~\eqref{eq:3-241} clearly allows at least one real root $< 0$; denote this root by $-d$. When equation~\eqref{eq:3-241} allows a real negative root (such as $-d$), we can factorize it in the form
\begin{equation}\label{eq:3-243}
(\sigma_1 + d)(\sigma_1^2 + 2b\sigma_1 + c)(\sigma_1 + d) = 0.
\end{equation}

By comparison with equation~\eqref{eq:3-241}, we obtain the relations
\begin{equation}\label{eq:3-244}
B = 2b + d, \quad C = 2bd + c, \quad \text{and} \quad D = cd.
\end{equation}

From these relations it follows that
\begin{equation}\label{eq:3-245}
BC - D = 2b(2bd + c + d^2) = 2b(C + d^2);
\end{equation}
therefore,
\begin{equation}\label{eq:3-246}
b = \frac{BC - D}{2(C + d^2)} \quad \text{and} \quad c = \frac{D}{d}.
\end{equation}

As we have noted, when $R_1 = 0$, $BC - D > 0$ and $D > 0$; therefore
\begin{equation}\label{eq:3-247}
b > 0 \quad \text{and} \quad c > 0 \quad (\text{when } R_1 \to 0).
\end{equation}

The roots of equation~\eqref{eq:3-243}, besides $-d$, are
\begin{equation}\label{eq:3-248a}
\sigma_1 = -b \pm \sqrt{b^2 - c};
\end{equation}
by~\eqref{eq:3-247}, these two roots have also negative real parts. The state $R_1 = 0$ is, therefore, \emph{absolutely} stable since $\mathrm{re}(\sigma_1) < 0$ for all three roots.

Consider what happens when $R_1$ increases: $D$ and $BC - D$ decrease and one of them may become zero. It is evident that if a real root of~\eqref{eq:3-243} becomes zero, $D \to 0$, since either $c \to 0$ or $d \to 0$; while if the real part of a pair of complex roots becomes zero $b \to 0$, i.e.\ $BC - D \to 0$ (since $C > 0$, $C + d^2$ is necessarily positive). Therefore, as $R_1$ increases, \emph{either $D \to 0$ (in which case a real root tends to zero) or $BC - D \to 0$ (in which case the real part of a pair of complex roots tends to zero)}; and depending on whichever happens first, we shall have the onset of instability as stationary convection or as overstable oscillations. But this is exactly the premise on which our earlier discussion was based; it is now justified.

%% --- Section 30 ---
\section{On a method of discriminating the character of the marginal state. A general variational principle}\label{sec:3-30}

The discussion in the preceding section was restricted to the case of two free boundaries for which it was possible to obtain the required characteristic equation in an explicit form. We now return to the general case when this is not possible.

Replacing $\sigma$ by $i\sigma$ in equations~\eqref{eq:3-93}--\eqref{eq:3-95} and letting
\begin{equation}\label{eq:3-248}
F = (D^2 - a^2)(D^2 - a^2 - i\sigma)W - \bigg(\frac{2\Omega}{\nu}\,d^2\bigg) DZ,
\end{equation}
we have (cf.\ equations~\eqref{eq:3-105} and~\eqref{eq:3-106})
\begin{equation}\label{eq:3-249}
(D^2 - a^2 - i\sigma)Z = -\bigg(\frac{2\Omega}{\nu}\,d\bigg) DW
\end{equation}
and
\begin{equation}\label{eq:3-250}
(D^2 - a^2 - ip\sigma)F = -Ra^2 W.
\end{equation}

Equations~\eqref{eq:3-248} and~\eqref{eq:3-249} can be combined to give
\begin{equation}\label{eq:3-251}
[(D^2 - a^2 - i\sigma)^2(D^2 - a^2) + \mathscr{T} D^2] W = (D^2 - a^2 - i\sigma)F.
\end{equation}

Solutions of equations~\eqref{eq:3-248}--\eqref{eq:3-250} must be sought which satisfy the boundary conditions
\begin{equation}\label{eq:3-252}
\begin{aligned}
&W = F = 0 \quad \text{for} \quad z = 0 \;\text{and}\; 1, \\
\text{and}\quad &\textit{either}\; DW = 0 \;\text{and}\; Z = 0 \quad \text{(on a rigid surface),} \\
&\textit{or}\; D^2 W = 0 \;\text{and}\; DZ = 0 \quad \text{(on a free surface).}
\end{aligned}
\end{equation}

There are eight boundary conditions and the requirement that a solution of equations~\eqref{eq:3-248}--\eqref{eq:3-250} satisfying these conditions will determine for a given $a^2$ and $i\sigma$ (which can be real or complex) a sequence of possible values for $R$. These characteristic values of $R$ will in general be complex; and the requirement that $R$ be real implies a relation between the real and imaginary parts of $\sigma$. We are not, however, interested in the relationship which may exist between the real and imaginary parts of $\sigma$ and the other parameters of the problem. We are interested only in the marginal state and its characterization. From the discussion of the case of two free boundaries in \S~29, it is clear that it will suffice for our present purposes to determine the critical Rayleigh number for the onset of instability as stationary convection, and as overstability, for a given Taylor number; that which gives the lower Rayleigh number is the one that will be manifested first as the Rayleigh number is gradually increased, a disturbance in the horizontal plane characterized by the wave number $a$ (in the unit $1/d$) first becomes unstable by overstability when it reaches the value $R_0(a^2)$; and $\sigma_0(a^2)$ is the frequency (in the unit $\nu/d^2$) of the oscillations which are set up in the marginal state. To determine the critical Rayleigh number for the onset of overstability, we must determine the minimum of the function $R_0(a^2)$. We should then compare the minimum so obtained with the critical Rayleigh number for the onset of stationary convection; and that which is lower is the one that will prevail and be manifested at the onset of instability.

It should be apparent from the foregoing algorithm that an exact solution of the double characteristic value problem will be difficult if not impracticable. However, a variational method can be devised which makes the problem feasible.

\subsection*{(a) The variational principle}

Equations~\eqref{eq:3-248}--\eqref{eq:3-250} can be treated in the same manner as equations~\eqref{eq:3-104}--\eqref{eq:3-106} were treated in \S~26\,$(a)$. Thus, by multiplying equation~\eqref{eq:3-250}

\noindent by $F$ and integrating over the range of $z$, we obtain after an integration by parts the equation
\begin{equation}\label{eq:3-253}
\int_0^1 [(DF)^2 + (a^2 + ip\sigma)F^2]\,dz = Ra^2 \int_0^1 WF\,dz.
\end{equation}

The right-hand side of this equation requires us to consider
\begin{equation}\label{eq:3-254}
\int_0^1 WF\,dz = \int_0^1 W\bigg[(D^2 - a^2)^2 W - i\sigma(D^2 - a^2)W - \frac{2\Omega}{\nu}\,d^2\,DZ\bigg]\,dz.
\end{equation}

After several integrations by parts we find that
\begin{equation}\label{eq:3-255}
\int_0^1 WF\,dz = \int_0^1 \{[(D^2 - a^2)W]^2 + i\sigma[(DW)^2 + a^2 W^2]\}\,dz + \bigg(\frac{2\Omega}{\nu}\,d^2\bigg)\int_0^1 Z\,DW\,dz,
\end{equation}
the integrated parts always vanishing on account of the boundary conditions. Now making use of equation~\eqref{eq:3-249}, we find after further integrations by parts that
\begin{equation}\label{eq:3-256}
\int_0^1 WF\,dz = \int_0^1 \{[(D^2 - a^2)W]^2 + d^2 [(DZ)^2 + a^2 Z^2]\}\,dz + i\sigma \int_0^1 \{(DW)^2 + a^2 W^2 + d^2 Z^2\}\,dz.
\end{equation}

Thus, the result of multiplying equation~\eqref{eq:3-250} by $F$ and integrating over the range of $z$ is
\begin{equation}\label{eq:3-257}
R = \frac{\displaystyle\int_0^1 [(DF)^2 + a^2 F^2 + ip\sigma F^2]\,dz}{a^2\;\displaystyle\bigg\{\int_0^1 \{[(D^2 - a^2)W]^2 + d^2 [(DZ)^2 + a^2 Z^2]\}\,dz + i\sigma \displaystyle\int_0^1 \{(DW)^2 + a^2 W^2 + d^2 Z^2\}\,dz\bigg\}}
= -\frac{I_1}{a^2 I_2} \;\text{(say)}.
\end{equation}

Now it can be shown (exactly as in \S~26\,$(a)$) that the variation $\delta R$ in $R$ given by equation~\eqref{eq:3-257} due to variations $\delta W$ and $\delta Z$ in $W$ and $Z$ compatible only with the boundary conditions on $W$, $Z$, and $F$ is given by
\begin{equation}\label{eq:3-258}
\delta R = -\frac{2}{a^2 I_2}\int_0^1 \delta F\{(D^2 - a^2 - ip\sigma)F + Ra^2 W\}\,dz.
\end{equation}

Consequently, \emph{if $\delta R = 0$ for all small arbitrary variations $\delta F$, then}
\begin{equation}\label{eq:3-259}
(D^2 - a^2 - ip\sigma)F = -Ra^2 W
\end{equation}
\emph{and conversely}. In other words, the characteristic values of $R$, for assigned $\sigma$, $a^2$, and $T$, have an extremal character; but unlike the case considered in \S~26\,$(a)$, they do not have a minimal character: indeed, they could not since they can be complex.

%% --- Section 31 ---
\section{The onset of convection as overstability: the solution for other boundary conditions}\label{sec:3-31}

The solution of the double characteristic value problem formulated in \S~30 can be obtained by the variational method by a procedure exactly analogous to that used in \S\S~27\,$(b)$ and~$(c)$. Thus, for obtaining the solution for the case of two rigid boundaries, we expand $F$ in a cosine series, as\footnote{The origin of $z$ has now been displaced so that its limits are $\pm\frac{1}{2}$.}
\begin{equation}\label{eq:3-260}
F = \sum_m A_m \cos\{(2m+1)\pi z\},
\end{equation}
and express $W$ and $Z$ as sums in the manner,
\begin{equation}\label{eq:3-261}
W = \sum A_m W_m \quad \text{and} \quad Z = \sum A_m Z_m.
\end{equation}

In view of equation~\eqref{eq:3-251}, $W_m$ is now a solution of the equation
\begin{equation}\label{eq:3-262}
[(D^2 - a^2 - i\sigma)^2(D^2 - a^2) + \mathscr{T} D^2] W_m = -c_{2m+1} \cos\{(2m+1)\pi z\},
\end{equation}
where
\begin{equation}\label{eq:3-263}
c_{2m+1} = (2m+1)^2\pi^2 + a^2 + i\sigma.
\end{equation}

The solution of equation~\eqref{eq:3-262} appropriate to the problem on hand is
\begin{equation}\label{eq:3-264}
W_m = c_{2m+1}\,\gamma_{2m+1}\cos\{(2m+1)\pi z\} + \sum_{j=1}^{3} B_j^{(m)} \cosh q_j\,z,
\end{equation}
where
\begin{equation}\label{eq:3-265}
\frac{1}{\gamma_{2m+1}} = c_{2m+1}^2[(2m+1)^2\pi^2 + a^2] + (2m+1)^2\pi^2\mathscr{T},
\end{equation}
the $B_j^{(m)}$ ($j = 1, 2, 3$) are constants of integration; and $q_j^2$ ($j = 1, 2, 3$) are the roots of the cubic equation,
\begin{equation}\label{eq:3-266}
(q^2 - a^2 - i\sigma)^2(q^2 - a^2) + \mathscr{T} q^2 = 0.
\end{equation}

The corresponding solution for $Z_m$ is
\begin{equation}\label{eq:3-267}
Z_m = -\frac{1}{\sqrt{\mathscr{T}}}\bigg\{(2m+1)\pi\gamma_{2m+1}\sin[(2m+1)\pi z] + \sum_{j=1}^{3}B_j^{(m)}\frac{q_j}{\chi_j}\sinh q_j\,z\bigg\},
\end{equation}
where
\begin{equation}\label{eq:3-268}
\chi_j = q_j^2 - a^2 - i\sigma \quad (j = 1, 2, 3).
\end{equation}

Comparing the solutions~\eqref{eq:3-264} and~\eqref{eq:3-267} with the solutions~\eqref{eq:3-146} and~\eqref{eq:3-161} obtained in \S~27\,$(a)$, we observe that they are identical except for

\noindent the different definitions of $q_j$, $\chi_j$, $c_{2m+1}$, and $\gamma_{2m+1}$. With these same redefinitions of the constants, the solutions for $B_j^{(m)}$ given in equations~\eqref{eq:3-154} and~\eqref{eq:3-155} apply equally well to the present case.

Substituting the solution for $W$ we have obtained in equation~\eqref{eq:3-250}, we now have
\begin{equation}\label{eq:3-269}
\begin{split}
&\sum A_m \bigg\{c_{2m+1}\gamma_{2m+1}\cos[(2m+1)\pi z] + \sum_{j=1}^3 B_j^{(m)}\cosh q_j\,z\bigg\} \\
&\qquad = Ra^2\sum A_m \bigg\{c_{2m+1}\gamma_{2m+1}\cos[(2m+1)\pi z] + \sum_{j=1}^3 B_j^{(m)}\cosh q_j\,z\bigg\}.
\end{split}
\end{equation}

Equation~\eqref{eq:3-269} is identical in form with equation~\eqref{eq:3-156}. The analysis following equation~\eqref{eq:3-156} up to, and inclusive of, equation~\eqref{eq:3-161} now applies.\footnote{Equations following equation~\eqref{eq:3-161} do not apply since these depend on the particular definitions of the roots $q_j$, etc.} In particular, we have the secular determinant:
\begin{equation}\label{eq:3-270}
\bigg\|\frac{1}{Ra^2}\,\frac{(2n+1)^2\pi^2 + a^2 + ip\sigma}{c_{2m+1}\gamma_{2m+1}}\,\delta_{mn} - \langle m|n\rangle\bigg\| = 0,
\end{equation}
where the matrix element $\langle m|n\rangle$ is given by equation~\eqref{eq:3-161}. The various quantities (such as $q_j$, $\chi_j$, $c_{2n+1}$, and $\gamma_{2m+1}$) appearing in equation~\eqref{eq:3-161} have the meanings now assigned to them.

The matrix $\langle m|n\rangle$ is symmetrical in $n$ and $m$. But the elements are complex, the symmetry does not ensure the reality of its characteristic roots; indeed they will in general be complex.

In the first approximation in which we retain only the $(0, 0)$-element of the secular matrix, the characteristic equation for $R$ becomes
\begin{equation}\label{eq:3-271}
R = \frac{1}{a^2}\,\frac{\pi^2 + a^2 + ip\sigma}{c_1\,\gamma_1\cdot\langle 0|0\rangle},
\end{equation}
where (cf.\ equation~\eqref{eq:3-161})
\begin{equation}\label{eq:3-272}
\begin{split}
\langle 0|0\rangle &= 2a^2_4 \Delta^{-1} \coth\tfrac{1}{2}q_1\bigg\{\tfrac{1}{2}(x_2 - x_1)(c_1 + x_3)^2 q_2\tanh\tfrac{1}{2}q_1 + {} \\
&\qquad + \tfrac{1}{2}(x_3 - x_2)(c_1 + x_2)^2 q_3\tanh\tfrac{1}{2}q_3 + \tfrac{1}{2}(x_3 - x_1)(c_1 + x_2)^2 q_2\tanh\tfrac{1}{2}q_2\bigg\}.
\end{split}
\end{equation}

It may be recalled that $\Delta$ in~\eqref{eq:3-272} is defined as in equation~\eqref{eq:3-155} but with the present meanings for the various quantities.

From our previous experience with variational methods, like the present, we may be confident that already the first approximation should give the required characteristic values to an accuracy of 1 or 2 per cent.

We shall now present the results of certain calculations based on equations~\eqref{eq:3-271} and~\eqref{eq:3-272}. The value
\begin{equation}\label{eq:3-273}
p = 0{\cdot}025
\end{equation}
was used in these calculations; this is approximately the value of $p$ for mercury at ordinary room temperatures. Some details regarding the method by which the $(R_c, T)$-relation was derived may be given.

For a chosen value of $a$, $R$ was evaluated in accordance with equations~\eqref{eq:3-271} and~\eqref{eq:3-272} for various assigned values of $\sigma$; and the value of $\sigma$ for which $R$ is real was deduced by interpolation. Thus for $T = 10^{10}$ and $a = 19{\cdot}8$ it was found that:
\begin{equation}\label{eq:3-274}
\begin{aligned}
R &= 1{\cdot}933 \times 10^6 - 5{\cdot}882 \times 10^4 i \quad \text{for } \sigma = 1{\cdot}420 \times 10^4; \\
R &= 1{\cdot}322 \times 10^6 + 9{\cdot}510 \times 10^4 i \quad \text{for } \sigma = 1{\cdot}500 \times 10^5; \\
R &= 1{\cdot}414 \times 10^6 + 0{\cdot}329 \times 10^4 i \quad \text{for } \sigma = 1{\cdot}489 \times 10^4.
\end{aligned}
\end{equation}

\begin{table}[ht]
\centering
\caption{Critical Rayleigh numbers and related constants for the onset of overstability in case both bounding surfaces are rigid and $p = 0{\cdot}025$}
\label{tab:XII}
\small
\begin{tabular}{ccccc}
\hline
$T$ & $a_c$ & $\sigma$ & $R_c$ & $p|\sigma|$ \\
\hline
$10^4$ & $3{\cdot}08$ & $4{\cdot}45 \times 10^1$ & $4{\cdot}39 \times 10^3$ & $0{\cdot}8692$ \\
$10^5$ & $4{\cdot}09$ & $5{\cdot}82 \times 10^2$ & $9{\cdot}31 \times 10^3$ & $1{\cdot}046$ \\
$5 \times 10^5$ & $8{\cdot}20$ & $2{\cdot}43 \times 10^3$ & $6{\cdot}39 \times 10^4$ & $0{\cdot}8862$ \\
$2 \times 10^6$ & $10{\cdot}28$ & $3{\cdot}90 \times 10^3$ & $1{\cdot}38 \times 10^5$ & $0{\cdot}5541$ \\
$10^7$ & $13{\cdot}46$ & $6{\cdot}81 \times 10^3$ & $3{\cdot}54 \times 10^5$ & $0{\cdot}4310$ \\
$10^{10}$ & $19{\cdot}7$ & $1{\cdot}29 \times 10^4$ & $1{\cdot}41 \times 10^6$ & $0{\cdot}2992$ \\
$10^{11}$ & $28{\cdot}75$ & $3{\cdot}47 \times 10^4$ & $5{\cdot}83 \times 10^6$ & $0{\cdot}2086$ \\
$10^{12}$ & $41{\cdot}7$ & $7{\cdot}18 \times 10^4$ & $2{\cdot}44 \times 10^7$ & $0{\cdot}1435$ \\
\hline
\end{tabular}
\end{table}

From these values, it can be estimated that
\begin{equation}\label{eq:3-275}
R = 1{\cdot}418 \times 10^6 \quad \text{for} \quad \sigma = 1{\cdot}4886 \times 10^4.
\end{equation}

By repeating such calculations for other $a$'s, the minimum of $R$ as a function of $a$ may be determined. Thus, in the example considered, it was found that:
\begin{equation}\label{eq:3-276}
\begin{aligned}
&\text{for } a = 19{\cdot}6, \quad \sigma = 1{\cdot}503 \times 10^4, \quad R = 1{\cdot}41765 \times 10^6; \\
&\text{for } a = 19{\cdot}7, \quad \sigma = 1{\cdot}496 \times 10^4, \quad R = 1{\cdot}4175 \times 10^6; \\
&\text{for } a = 19{\cdot}8, \quad \sigma = 1{\cdot}489 \times 10^4, \quad R = 1{\cdot}4177 \times 10^6.
\end{aligned}
\end{equation}

Consequently, it may be concluded that for $T = 10^{10}$, the critical Rayleigh number for the onset of overstability is $1{\cdot}4175 \times 10^6$ when $a = 19{\cdot}7$ and $\sigma = 1{\cdot}496 \times 10^4$.

In Table~XII the results of such calculations are summarized. Similar, but less extensive, calculations were undertaken for the case when one of the bounding surfaces is free and the other is rigid. The results of these calculations are given in Table~XIII.

\begin{table}[ht]
\centering
\caption{Critical Rayleigh numbers and related constants for the onset of overstability in case one bounding surface is free and the other is rigid and $p = 0{\cdot}025$}
\label{tab:XIII}
\small
\begin{tabular}{ccccc}
\hline
$T$ & $a_c$ & $\sigma$ & $R_c$ & $p|\sigma|$ \\
\hline
$10^5$ & $5{\cdot}85$ & $1{\cdot}44 \times 10^3$ & $1{\cdot}71 \times 10^4$ & $0{\cdot}9125$ \\
$3 \times 10^5$ & $15{\cdot}38$ & $1{\cdot}64 \times 10^4$ & $4{\cdot}81 \times 10^4$ & $0{\cdot}3601$ \\
$10^6$ & $0{\cdot}5$ & $7{\cdot}46 \times 10^4$ & $4{\cdot}87 \times 10^5$ & $0{\cdot}1492$ \\
\hline
\end{tabular}
\end{table}

The results on the critical Rayleigh number, for the onset of instability as stationary convection and as overstable oscillations for all three sets of boundary conditions, are all included in Figs.~21 and~22 (p.~96).

%% --- Section 32 ---
\section{The case $\mathfrak{p} = 0$}\label{sec:3-32}

The case when the Prandtl number $p$ tends to zero is singular with respect to the onset of overstable oscillations: thus, the asymptotic relations~\eqref{eq:3-234}, \eqref{eq:3-235}, and~\eqref{eq:3-237} are valid for $p^2 T \to \infty$, and they cannot be used when $p = 0$. The case $p = 0$ should be treated separately.

Returning then to equation~\eqref{eq:3-222} and setting $p = 0$, we have\footnote{We are restricting our present considerations to the case of two free boundaries.}
\begin{equation}\label{eq:3-277}
R_c^{(o)} = 2\pi^4\,\frac{(1+x)^3}{x};
\end{equation}
apart from the factor 2, this is the same formula which obtains in the absence of rotation (cf.\ equation~II\;\eqref{eq:3-193}). The minimum Rayleigh number occurs for $x = \frac{1}{2}$ when
\begin{equation}\label{eq:3-278}
R_c^{(o)} = 13{\cdot}5\pi^4 = 1315.
\end{equation}

The frequency of the overstable oscillations, when instability occurs, can be obtained from equation~\eqref{eq:3-221} by putting $p = 0$ and $x = \frac{1}{2}$; we get
\begin{equation}\label{eq:3-279}
\sigma_1^2 = \tfrac{4}{9}T_1 - 2{\cdot}25.
\end{equation}

Ignoring $2{\cdot}25$ in this equation and making use of the relation~\eqref{eq:3-236}, we have
\begin{equation}\label{eq:3-280}
p = (\tfrac{2}{3})^{\frac{1}{2}} 2\Omega,
\end{equation}
where $p$ is now the (circular) frequency in seconds.

The frequency given by equation~\eqref{eq:3-280} is different from the natural frequency of the rotating fluid only by the factor $(2/3)^{\frac{1}{2}}$. It would, therefore, be physically correct to say that \emph{in the limit $p = 0$, instability sets in by exciting the natural modes of oscillation of the rotating fluid}.

%% --- Section 33 ---
\section{Thermodynamic significance of the variational principles}\label{sec:3-33}

We shall now show that as in the case of the simple B\'enard problem (\S~14), the variational principles established in \S\S~26\,$(a)$ and~30 have similar thermodynamic meanings.

Consider the average rate of viscous dissipation of energy in a unit column of the fluid and the average rate at which the buoyancy force, $g\alpha\Theta$ ($= g\alpha\beta$) releases energy in the same unit column. These are given by (equations~II~\eqref{eq:3-175} and~\eqref{eq:3-178})
\begin{equation}\label{eq:3-281}
\epsilon_v = -\frac{p\nu}{d^2}\int_0^1\big\{\langle u(D^2 - a^2)u\rangle + \langle u(D^2 - a^2)u\rangle + \langle v(D^2 - a^2)v\rangle\big\}\,dz
\end{equation}
and
\begin{equation}\label{eq:3-282}
\epsilon_g = \rho g\alpha\int_0^1 \langle\delta w\rangle\,dz,
\end{equation}
where angular brackets signify that the quantity enclosed is averaged over the horizontal plane.

Without loss of generality for our present purposes, we may suppose that the expressions for $u$, $v$, and $w$ are those given in equations~\eqref{eq:3-186} and~\eqref{eq:3-187}. According to these equations, we have, for example,
\begin{equation}\label{eq:3-283}
\int_0^1 \langle u(D^2 - a^2)u\rangle\,dz = -\frac{1}{4a^4}\int_0^1 \{a_x^2\,DW(D^2 - a^2)\,DW + a_y^2 d^2 Z(D^2 - a^2)Z\}\,dz.
\end{equation}

It should be noted that the cross terms involving the products
\[
DW(D^2-a^2)Z \quad \text{and} \quad Z(D^2-a^2)W
\]
do not make any contributions to the foregoing expression; this is due to the quite general circumstance that the waves in the horizontal plane associated with $DW$ and $Z$ in the solution for $\mathbf{u}_\perp$ are exactly out of phase.

We have an expression similar to~\eqref{eq:3-283} for the term in $v$ in equation~\eqref{eq:3-281}; also,
\begin{equation}\label{eq:3-284}
\int_0^1 \langle w(D^2 - a^2)w\rangle\,dz = \pm\tfrac{1}{4}\int_0^1 W(D^2 - a^2)W\,dz.
\end{equation}

Combining these results, we have (cf.\ equation~II~\eqref{eq:3-178})
\begin{equation}\label{eq:3-285}
\epsilon_v = -\frac{p\nu}{4d^2 a^2}\int_0^1 \{a^2 W(D^2 - a^2)W + DW(D^2 - a^2)DW + d^2 Z(D^2 - a^2)Z\}\,dz.
\end{equation}

After integrating by parts the second and the third terms on the right-hand side, we find
\begin{equation}\label{eq:3-286}
\epsilon_v = \frac{p\nu}{4d^2 a^2}\int_0^1 \{[(D^2 - a^2)W]^2 + d^2 [(DZ)^2 + a^2 Z^2]\}\,dz.
\end{equation}

In the further reduction of the expression for $\epsilon_g$, we must distinguish the stationary and the oscillatory cases.

\subsection*{(a) The case when the marginal state is stationary}

In this case, the equation relating $\Theta$ and $w$ is (cf.\ equation~\eqref{eq:3-86})
\begin{equation}\label{eq:3-287}
w = -\frac{\kappa}{\beta d^2}(D^2 - a^2)\Theta.
\end{equation}

Accordingly,
\begin{equation}\label{eq:3-288}
\epsilon_g = -\frac{\rho g\alpha\kappa}{4\beta d^2}\int_0^1 \langle \Theta(D^2 - a^2)\Theta\rangle\,dz = \frac{\rho g\alpha\kappa}{4\beta d^2}\int_0^1 \Theta(D^2 - a^2)\Theta\,dz.
\end{equation}

After an integration by parts, we have
\begin{equation}\label{eq:3-289}
\epsilon_g = \frac{\rho g\alpha\kappa}{4\beta d^2}\int_0^1 [(D\Theta)^2 + a^2\Theta^2]\,dz.
\end{equation}

On the other hand, according to equations~\eqref{eq:3-98} and~\eqref{eq:3-104},
\begin{equation}\label{eq:3-290}
\Theta = -\frac{\nu}{g\alpha d^2 a^2}\,F.
\end{equation}

In terms of $F$ the expression for $\epsilon_g$ becomes
\begin{equation}\label{eq:3-291}
\epsilon_g = \frac{p\kappa^2}{4g\alpha\beta d^4 a^2}\int_0^1 [(DF)^2 + a^2 F^2]\,dz.
\end{equation}

In a steady state, the kinetic energy dissipated by viscosity must be balanced by the internal energy released by the buoyancy force. We must, therefore, require
\begin{equation}\label{eq:3-292}
\epsilon_v = \epsilon_g.
\end{equation}

Equations~\eqref{eq:3-286} and~\eqref{eq:3-291} now give
\begin{equation}\label{eq:3-293}
R = \frac{g\alpha\beta d^4}{\kappa\nu} = \frac{\displaystyle\int_0^1 [(DF)^2 + a^2 F^2]\,dz}{a^2\;\displaystyle\int_0^1 \{[(D^2 - a^2)W]^2 + d^2 [(DZ)^2 + a^2 Z^2]\}\,dz},
\end{equation}

\noindent But this is exactly the same expression for $R$ which, according to the variational principle of \S~26\,$(a)$, attains its absolute minimum for the marginal state.

\subsection*{(b) The case when the marginal state is oscillatory}

The question now arises: how should one modify under circumstances of overstability, the principle which equates the rate of irreversible dissipation of energy with the rate of liberation of the thermodynamically available energy by the buoyancy force, as a criterion for marginal stability? It would appear natural that we generalize the principle along the following lines.

If the motion is periodic in time, then the kinetic energy of the fluid, as well as the release of the potential energy by the buoyancy force, will be subject to similar variations. Suppose that all quantities which describe the perturbation vary with a circular frequency $p$ so that all the amplitudes have a time-dependent factor $e^{ipt}$. Then
\begin{equation}\label{eq:3-294}
\mathbf{u}_i\frac{\partial \mathbf{u}_i}{\partial t} = ip\mathbf{u}_i^2.
\end{equation}

We must allow for this change in the kinetic energy (per unit mass) in writing an equation of energy balance. Thus, it would appear that the quantity we must equate to $\epsilon_v$ is
\begin{equation}\label{eq:3-295}
\epsilon_v + ipp\int_0^1 \langle\mathbf{u}_i^2\rangle\,dz.
\end{equation}

Reverting to $d$ as the unit of distance and defining $\sigma$ ($= pd^2/\nu$) as in equation~\eqref{eq:3-92}, we have to consider
\begin{equation}\label{eq:3-296}
\epsilon_v + i\sigma\frac{p\nu}{d^2}\int_0^1 \langle\mathbf{u}_i^2\rangle\,dz
\end{equation}
in place of $\epsilon_v$ for the stationary case. The criterion for the occurrence of an oscillatory marginal state should then be
\begin{equation}\label{eq:3-297}
\epsilon_v + i\sigma\frac{p\nu}{d^2}\int_0^1 \langle\mathbf{u}_i^2\rangle\,dz = \epsilon_g = g\alpha\rho\int_0^1 \langle\delta w\rangle\,dz,
\end{equation}
where in evaluating $\epsilon_g$ we must allow for the fact that $\Theta$ and $w$ are no longer stationary.

At first sight, the appearance of the imaginary $i$, and complex numbers generally, in the equation for the energy balance, may strike one as very odd. Its origin must be traced to the fact that the oscillations in the velocity and in the acceleration are out of phase. Consequently, the

\noindent excess (or defect) of energy dissipated during one phase of the cycle must be exactly compensated by a similar excess (or defect) of energy liberated in a synchronous manner. It is this need for synchronism which determines the period of oscillation, as well as the Rayleigh number, as characteristics of the marginal state. The mathematical counterpart of this is that the determination of $R$ and $\sigma$ is through the solution of a double characteristic value problem.

We shall now verify that the thermodynamic principle, as reformulated above, is in agreement with the variational principle of \S~30. The expression~\eqref{eq:3-286} for $\epsilon_v$ continues to be valid. And we must now evaluate the additional term in $\langle\mathbf{u}_i^2\rangle$ in equation~\eqref{eq:3-297}. Considering the contribution to the integral by the $x$-component of the velocity, we have (cf.\ equation~\eqref{eq:3-187})
\begin{equation}\label{eq:3-298}
\int_0^1 \langle u^2\rangle\,dz = \frac{1}{4a^4}\int_0^1 [a_x^2(DW)^2 + a_y^2 d^2 Z^2]\,dz.
\end{equation}

We have a similar contribution from the $y$-component of the velocity; and altogether we have
\begin{equation}\label{eq:3-299}
i\sigma\frac{p\nu}{d^2}\int_0^1 \langle\mathbf{u}_i^2\rangle\,dz = i\sigma\,\frac{p\nu}{4d^2 a^2}\int_0^1 [a^2 W^2 + a^2(DW)^2 + a^2 d^2 Z^2]\,dz.
\end{equation}

This must be added to the expression~\eqref{eq:3-286} giving $\epsilon_v$.

Turning to the evaluation of $\epsilon_g$, we must now use the relation (see equation~\eqref{eq:3-93})
\begin{equation}\label{eq:3-300}
(D^2 - a^2 - ip\sigma)\Theta = -\bigg(\frac{\beta\,d^2}{\kappa}\bigg)w
\end{equation}
for eliminating $w$. Accordingly, we now have
\begin{equation}\label{eq:3-301}
\epsilon_g = -\frac{g\alpha\rho\kappa}{4\beta d^2}\int_0^1 \Theta(D^2 - a^2 - ip\sigma)\Theta^*\,dz = \frac{g\alpha\rho\kappa}{4\beta d^2}\int_0^1 [(D\Theta)^2 + (a^2 + ip\sigma)\Theta^2]\,dz.
\end{equation}

The relation between $\Theta$ and $F$ is unaltered and is given, as before, by~\eqref{eq:3-290}. Thus,
\begin{equation}\label{eq:3-302}
\epsilon_g = \frac{p\kappa^2}{4g\alpha\beta d^4 a^2}\int_0^1 [(DF)^2 + (a^2 + ip\sigma)F^2]\,dz.
\end{equation}

Combining equations~\eqref{eq:3-286}, \eqref{eq:3-297}, \eqref{eq:3-299}, and~\eqref{eq:3-302}, we obtain
\begin{equation}\label{eq:3-303}
R = \frac{\displaystyle\int_0^1 [(DF)^2 + (a^2 + ip\sigma)F^2]\,dz}{a^2\;\displaystyle\bigg\{\int_0^1\{[(D^2 - a^2)W]^2 + d^2[(DZ)^2 + a^2 Z^2]\} + i\sigma[(DW)^2 + a^2 W^2 + d^2 Z^2]\}\,dz\bigg\}}
\end{equation}
and this is, indeed, the expression for $R$ which formed the basis of the variational principle of \S~30.

We are thus enabled to formulate the following general principle.

\emph{Thermal instability as stationary convection will set in at the minimum (adverse) temperature gradient which is necessary to maintain a balance between the rate of dissipation of energy by viscosity and the rate of liberation of the thermodynamically available energy by the buoyancy force acting on the fluid. Likewise, the onset of thermal instability will be as overstable oscillations if it is possible (at a lower adverse temperature gradient) to balance in a synchronous manner the periodically varying amounts of kinetic energy with similarly varying amounts of dissipation and liberation of energy.}

%% --- Section 34 ---
\section{The case when $\boldsymbol{\Omega}$ and $\mathbf{g}$ act in different directions}\label{sec:3-34}

We shall now briefly consider the case when $\boldsymbol{\Omega}$ and $\mathbf{g}$ act in different directions. For this purpose we must return to equations~\eqref{eq:3-84} and~\eqref{eq:3-85} in which the assumption that $\boldsymbol{\Omega}$ and $\mathbf{g}$ are parallel has not yet been made.

Let $\boldsymbol{\Omega}$ be inclined at an angle $\vartheta$ to the direction of the vertical. Also, let the direction of the $x$-axis be so chosen that $\boldsymbol{\Omega}$ lies in the $xz$-plane. Then
\begin{equation}\label{eq:3-304}
\boldsymbol{\lambda} = (0, 0, 1) \quad \text{and} \quad \boldsymbol{\Omega} = \Omega(\sin\vartheta, 0, \cos\vartheta),
\end{equation}
and equations~\eqref{eq:3-84} and~\eqref{eq:3-85} become
\begin{equation}\label{eq:3-305}
\frac{\partial}{\partial t}\nabla^2 w = \nu\nabla^4 w + 2\Omega\bigg(\cos\vartheta\,\frac{\partial \zeta}{\partial z} + \sin\vartheta\,\frac{\partial \zeta}{\partial x}\bigg) + g\alpha\nabla_1^2\Theta
\end{equation}
and
\begin{equation}\label{eq:3-306}
\frac{\partial}{\partial t}\nabla^2 w = g\alpha\bigg(\frac{\partial^2}{\partial x^2} + \frac{\partial^2}{\partial y^2}\bigg)\Theta + \nu\nabla^4 w - 2\Omega\bigg(\cos\vartheta\,\frac{\partial^2}{\partial z^2} + \sin\vartheta\,\frac{\partial}{\partial x}\frac{\partial}{\partial z}\bigg)\zeta.
\end{equation}

The remaining equation~\eqref{eq:3-86} is unaffected by this generalization.

If we seek solutions of equations~\eqref{eq:3-86}, \eqref{eq:3-305}, and~\eqref{eq:3-306} which are independent of $x$ and are of the forms
\begin{equation}\label{eq:3-307}
w = W(z)\cos\frac{a_y}{d}\,y, \quad \zeta = Z(z)\cos\frac{a_y}{d}\,y, \quad \text{and} \quad \Theta = \Theta(z)\cos\frac{a_y}{d}\,y,
\end{equation}
then equations~\eqref{eq:3-305} and~\eqref{eq:3-306} reduce to equations~\eqref{eq:3-87} and~\eqref{eq:3-88} with the only difference that $\Omega$ is replaced by $\Omega\cos\vartheta$. Consequently, if we restrict ourselves to the onset of instability as rolls in the $x$-direction, then all of the discussion of the preceding sections, with the exception of those parts of \S~28 which specifically deal with general cell patterns, will apply if we interpret $\Omega$ everywhere to mean the component of $\boldsymbol{\Omega}$ in the direction of $\mathbf{g}$.

The question, whether lower Rayleigh numbers can be attained for the

\noindent marginal states by considering more general patterns of motion than rolls, can be answered only by solving the necessary characteristic value problems in which $a_x$ and $a_y$ are explicitly retained; and minimizing the lowest characteristic values, for given $a_x$ and $a_y$, as functions of these two parameters. Such an analysis has not been carried out for the present problem; but it has been carried out for the closely related problem of the inhibition of thermal convection by an impressed magnetic field (Chapter~IV, \S~47). If the conclusions reached in that context can be extended to the present problem, then it would appear that the minimum Rayleigh numbers do indeed occur for the onset of instability as rolls. This apparently paradoxical conclusion is clarified in Chapter~IV, \S~47.

%% --- Section 35 ---
\section{Experiments on the onset of thermal instability in rotating fluids}\label{sec:3-35}

We have seen that rotation influences the onset of thermal instability
in many ways.  The principal effects are these.

\begin{enumerate}
\item[(1)] Rotation inhibits the onset of instability; the extent of the
  inhibition depends on the Taylor number $\mathscr{T}$ ($= 4\Omega^2 d^4/\nu^2$) and the
  Prandtl number $p$ ($= \nu/\kappa$); and as $\mathscr{T} \to \infty$, the Rayleigh number at which
  instability sets in, and the wave number $a$ (in units of $1/d$) with which it
  is associated, have the asymptotic behaviours
  \begin{equation}
    R \to \text{constant}\; \mathscr{T}^{\!\nicefrac{2}{3}}
    \quad \text{and} \quad
    a \to \text{constant}\; \mathscr{T}^{\!\nicefrac{1}{6}}.
    \label{eq:3-308}
  \end{equation}

\item[(2)] The onset of instability will be as stationary convection so long
  as $p$ exceeds a certain critical value $p^*$.  The precise value of $p^*$ depends
  on the nature of the bounding surfaces; but it can be inferred from the
  value ($= 0.6766$) determined for the case of two free boundaries that
  $p^* \sim 1$.  When $p > p^*$, the $(R_c, \mathscr{T})$-relation is independent of $p$.

\item[(3)] The onset of instability will be as overstable oscillations if $p < p^*$
  and the Taylor number exceeds a certain value $\mathscr{T}^{(0)}$ depending on $p$.  For
  $\mathscr{T} \leq \mathscr{T}^{(0)}$ the onset of instability will be as stationary convection.  For a
  given $p$ ($< p^*$) the asymptotic behaviours of $R$ and $a$ are again as in (1)
  above; but the constants of proportionality now depend on $p$.  Also, the
  gyration frequency $p$ of the overstable oscillations is essentially deter\-mined
  by $\Omega$ and the Taylor number; for $0 < p < p^*$,
  \begin{equation}
    \nu/\Omega \to \text{constant}\; \mathscr{T}^{-\nicefrac{1}{6}}
    \quad \text{as} \quad \mathscr{T} \to \infty,
    \label{eq:3-309}
  \end{equation}
  where the constant depends on $p$.
\end{enumerate}

From the foregoing summary of the principal effects to be expected,
it is clear that in selecting the liquids for the experiments it would be
preferable to choose two liquids, one with a Prandtl number substantially
larger than unity and the other with a Prandtl number substantially less
than unity.  Two such liquids are readily available: water with $p \sim 7{\cdot}5$
and mercury with $p = 0{\cdot}025$.  Experiments with these two liquids have
been carried out by Nakagawa and Frenzen, Fultz and Nakagawa, and
Goroff.  An account of these experiments will now be given.

%----------------------------------------------------------------------
\subsection{Experiments with water}
\label{sec:35a}
%----------------------------------------------------------------------

In these experiments, water was contained in Pyrex glass cylinders
of 12 inches outside diameter.  To the bottoms of the cylinders were
cemented electrically conducting Pyrex heating plates; they served as
the heating elements to which measured amounts of electric power could
be supplied.  The cylinders containing the fluid and the necessary
accessories were placed on a supporting steel ring mounted on a vertical
shaft.  The steel ring had provisions for levelling and for other adjust\-ments.
The fluid was set in rotation by driving the shaft with suitable
belts and motors.  In the experiments, layers of water of depths varying
from 3 to 17~cm and speeds of rotation up to 60~rev/min were used.  Taylor
numbers up to $10^8$ were reached in this manner.

The difference in the temperatures at two levels, one near the top and
one near the bottom surface of the fluid, was measured by copper--constantan
wire thermocouples.  By using several junctions in series as
a thermopile, it was possible to measure the difference in the average
temperatures at the two levels with an accuracy of $\pm 0{\cdot}01$~$^\circ$C; this
accuracy in the temperature measurements was found to be necessary
and sufficient.

Observations with the foregoing experimental arrangement were of
two kinds: visual and photographic observations made with the aid of a
rotoscope which enabled the direct detection of the fluid motions
relative to the rotating frame; and quantitative observations based on
the records of the difference of temperature between the top and bottom
surfaces made continuously during each experiment after varying
amounts of electric power had been turned on to the heating elements.

By visual observations with the rotoscope, it was possible to observe
the onset of cellular convection when the heating rate exceeded a certain
amount.  The motions could be made visible by scattering a small
amount of aluminium powder on the top surface.  The movements of the
aluminium particles could be photographed by a time-exposure streak
technique.  An example of such a photograph is shown in Fig.~30.  The
internal motions in the cells could also be followed with black ink and

\figurePlaceholder{30}{Convection cells which appear in water when
in rotation and heated from below.  The data to which
the illustration refers are: depth 18~cm, difference in
temperature $0{\cdot}7^\circ$; rate of rotation $5{\cdot}0$~rev/min; Taylor
number $1{\cdot}2 \times 10^8$.}

\figurePlaceholder{31}{Side view of ink rising from the bottom of a
cylinder of water in the ascending cores of cells.  The
data to which the illustration refers are: depth 18~cm,
difference in temperature $0{\cdot}5^\circ$; rate of rotation $10{\cdot}9$~rev/min;
Taylor number $5{\cdot}5 \times 10^9$.}

\noindent photographed.  Fig.~31 is an example of a photograph showing such
internal motions.

For a quantitative determination of the critical Rayleigh number for
the onset of instability, a careful study of the temperature records is
necessary.  In Fig.~32 examples of temperature records which were
obtained for varying heating rates are shown.  Fig.~32\,$a$ corresponds to

\figurePlaceholder{32}{Time records of the adverse temperature gradient for water for three
different rates of heating ($d = 3$~cm, $\Omega = 10$~rev/min) (from \emph{Proc.\ Roy.\ Soc.\
(London)} A, \textbf{231}, 519 (1955)).}

\noindent a case when the bottom plate was heated at a sufficiently low rate that
the steady difference in temperature attained was insufficient to induce
instability; Fig.~32\,$b$ is a typical record when the heating rate sufficed to
surpass by some margin the temperature gradient necessary for the
onset of instability; Fig.~32\,$c$ is an example when the heating rate was so
much in excess of what is necessary to cause instability that the conditions
were approaching turbulence.  The difference between the records, when
we are in the conductive r\'egime and when we have passed it, is sufficiently
striking that by examining them for various heating rates, the tempera\-ture
gradient at which conditions are just past marginal could be deter\-mined
with some precision.  Simultaneous visual observations confirmed
the criteria that were used.  In this manner, Nakagawa and Frenzen

\figurePlaceholder{33}{Comparison of the theoretical relation with the experimental results
on the onset of thermal instability in water in varying speeds of rotation: the
critical Rayleigh number $R_c$ is plotted against the Taylor number $\mathscr{T}$.  Experi\-ments
with different depths of water are distinguished: $\boxtimes$ 18~cm; $\diamondsuit$ 14~cm;
$\odot$ 10~cm; $\bigotimes$ 6~cm; $\times$ 17.3~cm; $\triangle$ 13.2~cm; $\blacksquare$ 9.3~cm; $\circ$ 6.3~cm; $\boxminus$ 3~cm;
$\bigcirc$ 2~cm.}

\noindent determined the critical Rayleigh number as a function of the Taylor
number.  Their results are shown in Fig.~33.  The theoretical $(R_c, \mathscr{T})$-relation
for one rigid and one free surface derived in \S\,27\,$(c)$ is shown in
the same plot; it will be observed that the theoretically predicted relation
is fully confirmed.

Estimates of the cell dimensions based on photographic observations
are also in accord with the theoretical expectations.

%----------------------------------------------------------------------
\subsection{Experiments with mercury}
\label{sec:35b}
%----------------------------------------------------------------------

The experiments of Fultz and Nakagawa with mercury were carried
out with essentially the same experimental arrangement as described
in \S\,(a) above.  But a number of precautions peculiar to mercury had to
be observed.  For example, Fultz and Nakagawa record that considerable
trouble was experienced with the oxidation of the mercury surface.  It
was necessary on this latter account to replace the air above the heated
surface by nitrogen.  This was accomplished by circulating nitrogen;
this also served the purpose of keeping the temperature of the top surface
constant.  Even with these precautions, it was found that some time after
the heating current had been turned on, a thin contaminant film was
formed on the surface.  The film was rigid enough to prevent motions at
the top surface.  The layer of mercury was thus effectively confined
between two rigid surfaces.

The temperature gradients at which instability occurred were again
determined by an examination of the temperature records.  This was
supplemented by making plots of the temperature gradients against the
heating rates and using the Schmidt--Milverton principle (Chapter~II,

\figurePlaceholder{34}{Time records of the adverse temperature gradient for mercury for three
different rates of heating ($d = 6$~cm, $\Omega = 15$~rev/min) (from \emph{Proc.\ Roy.\ Soc.\
(London)} A, \textbf{231}, 520 (1955)).}

\noindent \S\,18\,(b)).  Layers of mercury of depth up to 8~cm and speeds of rotation
up to 30~rev/min were used.  Taylor numbers in the range $10^6$--$10^{12}$ were
thus accessible to the experiments.

In Fig.~34 a series of temperature records similar to those in Fig.~32
for water are shown.  The difference in the two sets of temperature
records, after instability has set in, is very striking; in the experiments
with mercury, the temperature records exhibit very pronounced periodic
fluctuations.  Fig.~35, exhibiting records which show these periodic
fluctuations with uniformity over durations as long as half an hour and
more, leaves very little doubt that we are witnessing here the overstable
oscillations predicted by the theory.  Indeed, by counting the number

\figurePlaceholder{35}{Further time records of the adverse temperature gradient for mercury.}

\noindent The conditions to which they refer are given below:
\begin{center}
\begin{tabular}{lcc}
  & (a) & (b) \\
depth (cm) & 6 & 6 \\
$\Omega$ (sec$^{-1}$) & $1{\cdot}11$ & $3{\cdot}08$ \\
temperature difference (degrees) & $1{\cdot}54$ & $5{\cdot}05$ \\
Taylor number ($\mathscr{T}$) & $5 \times 10^9$ & $3{\cdot}9 \times 10^{10}$ \\
Rayleigh number ($R$) & $2{\cdot}64 \times 10^4$ & $8{\cdot}94 \times 10^4$ \\
period of oscillation (sec) & $19$--$20$ & $9$--$10$
\end{tabular}
\end{center}

\noindent of these fluctuations in known intervals of time, we can estimate the
basic periods of the oscillations.  The results of such estimates are shown
in Fig.~36 together with the theoretically expected relation.  It will be
seen that the agreement is satisfactory.

\figurePlaceholder{36}{A comparison of the observed periods of oscillation at marginal stability
with the theoretical periods.  The ordinate $\tau$ gives the period in units of $2\varpi$.  The curves
$aa$ and $bb$ are the theoretical relations for $p = 0{\cdot}025$ and for the case of two bounding
surfaces rigid ($aa$) and one bounding surface rigid and the other free ($bb$).}

\figurePlaceholder{37}{A summary of the theoretical and experimental results on the critical Rayleigh
numbers for the onset of instability in mercury.  The curves $aa$, $bb$, and $cc$ are the
theoretical $(R_c, \mathscr{T})$-relations for the three cases (i) both bounding surfaces rigid, (ii) one
bounding surface rigid and the other free, and (iii) both bounding surfaces free for
$p = 0{\cdot}025$.  The solid circles are the experimentally determined points for mercury
($p = 0{\cdot}025$).}

In Fig.~37 the results on the critical Rayleigh numbers are shown
together with the predicted relations for the three sets of boundary
conditions.  It is noteworthy that the agreement is best with the
theoretical relation for two rigid boundaries: as we have already
remarked, the experimental conditions do correspond to this case.

The experiments of Nakagawa confirm that the onset of thermal

\figurePlaceholder{38a}{The variation of the Nusselt number for increasing Rayleigh numbers
for the case when the Taylor number is $1{\cdot}12 \times 10^9$.  The sharp break at the
Rayleigh number $3{\cdot}42 \times 10^4$ occurs while the system is already overstable.
Notice that the Nusselt number has not increased substantially beyond 1 even
though overstability set in at $R \sim 2{\cdot}5 \times 10^4$; this indicates the relative in\-efficiency
of heat-transfer by overstable oscillations.}

\figurePlaceholder{38b}{The dependence of the Rayleigh number on
the Taylor number at which the second break occurs.
The full-line curve is the theoretical relation for the
onset of stationary convection ignoring overstability.}

\noindent instability in a horizontal layer of mercury in rotation occurs as over\-stable
oscillations at the predicted Rayleigh numbers $R_c^{(o)}$ and with the
predicted characteristic frequencies.  When the Rayleigh number
exceeds $R_c^{(o)}$, the overstable oscillations will become of finite amplitude
and non-linear effects, ignored in the linear stability theory, will become
effective.  At the same time, the mode with the real exponential time
dependence will continue to be damped (on the assumption that the
presence of the overstable oscillations does not seriously alter the static
initial state to which the entire theory refers).  When $R$ becomes equal
to $R_c^{(s)}$ ($> R_c^{(o)}$), this latter mode will begin to manifest itself.  Since $R_c^{(s)}$
is about twenty-six times $R_c^{(o)}$ (for $\mathscr{T} \to \infty$, when both bounding surfaces
are rigid, and $p = 0{\cdot}025$), it is clear that the onset of stationary convec\-tion
will occur superposed on overstable oscillations of quite finite
amplitudes.  However, on general grounds, one may expect that the
overstable oscillations are not very efficient in transporting heat; and
that on this account one may, nevertheless, be able to detect the onset
of stationary convection in a Schmidt--Milverton plot.  Some experi\-ments
of Dropkin and Globe seemed to indicate that a second break\footnote{The first break was so weak in these experiments that one could not be certain of
it in most cases.}
occur in a Schmidt--Milverton plot after a first break corresponding to
the onset of overstability.  But these experiments were not carried out
with sufficient precautions to decide whether or not the second break
occurs at the predicted Rayleigh number $R_c^{(s)}$.  More recent experiments
by Goroff carried out with the necessary care confirms that the second
break does occur at the predicted place.  In Fig.~38\,$a$ the results of his
experiments for $\mathscr{T} = 1{\cdot}12 \times 10^9$ are shown; in this plot of the Nusselt
number against the Rayleigh number, a break clearly occurs at
$R = 3{\cdot}45 \times 10^4$; and this is long after overstable oscillations had become
a permanent feature of the temperature records.  The value
$R = 3{\cdot}45 \times 10^4$ should be compared with the theoretical value,
$R_c^{(s)} = 3{\cdot}50 \times 10^4$ which the results of Table~VIII predict for the onset
of stationary convection for $\mathscr{T} = 1{\cdot}12 \times 10^9$.  In Fig.~38\,$b$ the values of
$R_c^{(s)}$ obtained for three different values of $\mathscr{T}$ (including the one to which
Fig.~38\,$a$ refers) are shown together with the theoretically expected
relation.  It would appear from these experiments that there can be
little doubt that the onset of stationary convection occurs at the pre\-dicted
Rayleigh number.  They further establish the relative inefficiency
of overstable oscillations for the transport of heat.

%% --- Bibliographical Notes ---
\section*{BIBLIOGRAPHICAL NOTES}
\addcontentsline{toc}{section}{Bibliographical Notes}

For a general account of the subject of this chapter see:

\begin{enumerate}
\item[1.] S.\ \textsc{Chandrasekhar}, `Thermal convection', \emph{Daedalus}, \textbf{86}, 323--39 (1957).
\end{enumerate}

\noindent \S\,22.  The fundamental theorem proved in this section is due to:
\begin{enumerate}
\item[2.] G.\ I.\ \textsc{Taylor}, `Experiments with rotating fluids', \emph{Proc.\ Roy.\ Soc.\ (London)}
  A, \textbf{100}, 114--21 (1921).
\end{enumerate}

\noindent The same theorem, though not in as explicit a form, was independently stated by:
\begin{enumerate}
\item[3.] J.\ \textsc{Proudman}, `On the motion of solids in a liquid possessing vorticity',
  \emph{Proc.\ Roy.\ Soc.\ (London)} A, \textbf{92}, 408--24 (1916).
\end{enumerate}

In Taylor's paper (reference~2), experiments are described which demonstrate
very effectively the implications of the theorem for the character of the fluid
motions which prevail under the circumstances.  Taylor's experiments have been
repeated by:
\begin{enumerate}
\item[4.] D.\ \textsc{Fultz}, `A survey of certain thermally and mechanically driven systems
  of meteorological interest', \emph{Proceedings of the First Symposium on the Use
  of Models in Geophysical Fluid Dynamics}, 27--63, edited by Robert B.\ Long,
  Johns Hopkins University, Baltimore, 1953.
\end{enumerate}

\noindent The following references to related papers by Taylor may be noted:
\begin{enumerate}
\item[5.] G.\ I.\ \textsc{Taylor}, `Motion of solids in fluids when the flow is not irrotational',
  \emph{Proc.\ Roy.\ Soc.\ (London)}, A, \textbf{93}, 99--113 (1917).
\item[6.] \rule{2em}{0.4pt}, `The motion of a sphere in a rotating liquid', ibid., \textbf{102}, 180--9
  (1922).
\item[7.] \rule{2em}{0.4pt}, `Experiments on the motion of solid bodies in rotating fluids', ibid.,
  \textbf{104}, 213--18 (1923).
\end{enumerate}

The dynamics of rotating fluids have many fascinating aspects of which the
Taylor--Proudman theorem is only one.  For a general account of these other
aspects see:
\begin{enumerate}
\item[8.] T.\ E.\ \textsc{Sterne}, `Rotating fluids', \emph{Surveys in Mechanics}, 139--61, edited by
  G.\ K.\ Batchelor and R.\ M.\ Davies, Cambridge Monographs on Mechanics
  and Applied Mathematics, Cambridge, England, 1956.
\end{enumerate}

\noindent See also:
\begin{enumerate}
\item[9.] G.\ W.\ \textsc{Morgan}, `A study of motions in a rotating liquid', \emph{Proc.\ Roy.\ Soc.\
  (London)} A, \textbf{206}, 108--30 (1951).
\item[10.] H.\ \textsc{G\"ortler}, `\"Uber eine Schwingungserscheinung in Fl\"ussigkeiten mit
  stabiler Dichteschichtung', \emph{Z.\ ang.\ Math.\ Mech.}\ \textbf{23}, 65--71 (1943).
\item[11.] \rule{2em}{0.4pt}, `Einige Bemerkungen \"uber Str\"omungen in rotierenden Fl\"ussigkeiten',
  ibid.\ \textbf{24}, 210--14 (1944).
\end{enumerate}

\noindent \S\S\,24--27.  The analyses in these sections are derived from:
\begin{enumerate}
\item[12.] S.\ \textsc{Chandrasekhar}, `The instability of a layer of fluid heated below and
  subject to Coriolis forces', \emph{Proc.\ Roy.\ Soc.\ (London)} A, \textbf{217}, 306--27 (1953).
\item[13.] \rule{2em}{0.4pt}\ and \textsc{Donna} D.\ \textsc{Elbert}, `The instability of a layer of fluid heated
  below and subject to Coriolis forces.\ II', ibid., \textbf{231}, 198--210 (1955).
\end{enumerate}

\noindent The method of solution described in \S\,27 (for cases $(b)$ and $(c)$) is, however, different
from that in reference~12; and the results included in Tables~VIII and IX (due
to Miss Donna Elbert) are based on the new formulae.

\noindent \S\,28.  The discussion of the cell patterns in this section follows:
\begin{enumerate}
\item[14.] G.\ \textsc{Veronis}, `Cellular convection with finite amplitude in a rotating fluid',
  \emph{J.\ Fluid Mech.}, \textbf{5}, 401--35 (1959).
\end{enumerate}

\noindent \S\,29.  The principal results of this section are contained in references~12 and~13.
But the arguments have been rearranged and made more explicit.  The particular
form of the discussion of the roots given in \S\,(a) is due to Dr.\ John Sykes.

\noindent \S\S\,30--32.  The analyses in these sections are again derived from references~12
and~13.

\noindent \S\,33.  See:
\begin{enumerate}
\item[15.] S.\ \textsc{Chandrasekhar}, `The thermodynamics of thermal instability in
  liquids', \emph{Max Planck Festschrift 1958}, 103--14, Veb Deutscher Verlag der
  Wissenschaften, Berlin, 1958.
\end{enumerate}

\noindent \S\,35.  The references for the experimental work described in this section are:
\begin{enumerate}
\item[16.] D.\ \textsc{Fultz}, Y.\ \textsc{Nakagawa}, and P.\ \textsc{Frenzen}, `An instance in thermal
  convection of Eddington's ``overstability''\,', \emph{Physical Rev.}, \textbf{94}, 1471--2
  (1954).
\item[17.] \rule{2em}{0.4pt}, `Experiments on overstable thermal convection in mercury',
  \emph{Proc.\ Roy.\ Soc.\ (London)} A, \textbf{231}, 211--25 (1955).
\item[18.] Y.\ \textsc{Nakagawa} and P.\ \textsc{Frenzen}, `A theoretical and experimental study of
  cellular convection in rotating fluids', \emph{Tellus}, \textbf{7}, 1--21 (1955).
\item[19.] D.\ \textsc{Dropkin} and S.\ \textsc{Globe}, `Effect of spin on natural convection in
  mercury heated from below', \emph{J.\ Appl.\ Phys.}, \textbf{30}, 84--89 (1959).
\item[20.] I.\ R.\ \textsc{Goroff}, `An experiment on heat transfer by overstable and ordinary
  convection', \emph{Proc.\ Roy.\ Soc.\ (London)} A, \textbf{254}, 537--41 (1960).
\end{enumerate}
